\subsection*{19AZ. Freeze-audit verdict on the machine-theoretic line}
In this inherited release-baseline block, the machine-theoretic line reaches the release-level freeze threshold used for the then-present public version. A local audit of Sections~19J--19AY shows that the computational route itself is already complete enough for the present Companion purpose, while the last overlay block from 19AI onward increasingly repackages nearby closure properties rather than opening substantially new machine content. For this reason that inherited line should now be read with a strict two-tier discipline. Theorems~19J--19AH remain the indispensable machine ladder: local operational semantics, conditional Zuse-style primitive core, reference-core closure, relative tape-family upgrade, the strict upper gap, the full machine-theoretic closure surface, the status decoder, and the present-line conditional/unconditional status surface. The local continuation block 19AH$\prime$--19AH$^{(22)}$ should be read only as a diagnostic refinement layer: it separates total and imbalance channels, turns that split into a control dictionary, sorts inherited continuation knots by burden type, and compresses the surviving OP~1 queue to the present optimum-pinning comparison-rigidity kernel. Section~19AI now functions as the bundled overlay-completion capstone: it certifies in compressed form that the machine line is conservative, faithful, exact, relatively complete, adequate, reconstructive, closed, idempotent, minimal, canonical, invariant, natural, universally characterized, rigid, terminal, and equivalent to the certified machine-bearing windows on the intended admissible domain. The audit verdict in that inherited release block is therefore to stop expansion there: future revisions should preferentially cite the bundled surfaces already obtained --- in particular 19AF, 19AH, and 19AI --- rather than append further parallel overlay adjectives unless a genuinely new machine-theoretic distinction is introduced.

\subsection*{19AZ$\prime$. Post-release status note after tool/QAM completion and local OP-kernel closure}
The public v3.28.0 release completed the formerly deferred tool/QAM/decoder layer and fixed the route-control infrastructure of the compressed computational line. The present post-release working line then used that infrastructure exactly as intended: it reopened the two parked residual burdens through the smallest licensed routes and followed them to local kernel closure. OP~1 was normalized first as a retained pinning-signature injectivity problem, then split into occupancy-side and completion-side duplicate-signature witnesses, and finally discharged under the standing retained no-second-carrier and strict-upper-gap exclusions. OP~5 Step~(B) was normalized as a carrier-upgraded routed-balance witness problem, split into carrier-stable versus obstructed routed returns, then into upstream-screened versus downstream carrier-grade obstruction, and finally discharged on the retained branch-shell corridor by the inherited shell-closure package.

The honest status after these steps is therefore no longer ``two parked residual burdens.'' It is: the completed tool/QAM/decoder infrastructure now stands as fixed route-control surface, and the live OP~1 and OP~5 Step~(B) kernels have both been discharged locally under the standing retained disciplines. What remains conditional is not the local kernel closure but the broader machine-theoretic reading beyond the certified retained corridors and routes. The governing public DOI in the chronology is \href{https://doi.org/10.5281/zenodo.19699827}{10.5281/zenodo.19699827} for the public v3.31.0 release.

\subsection*{19AZ$''$. Immediate normalized channel-sorted tool ledger}
\label{sec:immediate-normalized-tool-ledger}
The first concrete task of the post-freeze continuation line is not to invent a fourth proof language, but to normalize the already extracted tools into a channel-sorted ledger. On the present line this means three reusable packets: one for statements that are already total-channel stable, one for statements that compress the surviving imbalance-side burden, and one for statements that translate between both surfaces or between their QAM/matrix realizations. The point is discipline: later arguments should be able to cite the smallest packet that already matches their burden class.

\begin{definition}[Normalized channel-sorted tool packets]
\label{def:normalized-channel-sorted-tool-packets}
The immediate continuation line organizes the current extracted tools into three provisional packets.
\begin{enumerate}[leftmargin=2em]
  \item \textbf{Total-channel packet} \(\mathfrak T_{\mathrm{tot}}\). This packet contains the tools whose role is to keep the retained total-channel architecture fixed once admissible transport has already selected the active branch. Its current generators are the transport-first tool, the readout-after-transport tool, the admissible-readout uniqueness tool, and the global pass/transfer/diagnostic-subordination interface from Propositions~\ref{prop:global-pass-tool}--\ref{prop:global-diagnostic-subordination-tool}.

  \item \textbf{Imbalance packet} \(\mathfrak T_{\mathrm{imb}}\). This packet contains the tools that compress the surviving OP~1 burden after the total-channel surface has been fixed. Its current generators are the frame-transport split certificate (Proposition~\ref{prop:op1-frame-transport-split}), the finite four-shell certificate (Proposition~\ref{prop:op1-four-shell-certificate}), the two-pass half-burden certificate (Proposition~\ref{prop:op1-two-pass-half-burden}), the reciprocal pass-symmetry certificate (Proposition~\ref{prop:op1-reciprocal-pass-symmetry}), the same-corridor reversal implication (Proposition~\ref{prop:op1-same-corridor-implies-reciprocity}), the unit-loop half-burden certificate (Proposition~\ref{prop:op1-unit-loop-half-burden}), the one-first-zero loop-budget bridge (Proposition~\ref{prop:op1-one-first-zero-loop-budget}), the admissible-window uniqueness bridge (Proposition~\ref{prop:op1-admissible-window-uniqueness}), the retained-window-rigidity / no-second-carrier / strict-upper-gap bridge (Proposition~\ref{prop:op1-rigidity-gap-uniqueness-bridge}), and the optimum-pinning comparison-rigidity bridge (Proposition~\ref{prop:op1-optimum-pinning-implies-rigidity}).

  \item \textbf{Bridge packet} \(\mathfrak T_{\mathrm{br}}\). This packet contains the tools whose job is to move between total and imbalance language, or between channel language and its shell/menu/matrix realizations. Its current generators are the shared matrix-leakage tool, the cryptanalytic translation tool, the shared phase-coupling tool, and on the QAM side the shell-stable readout tool, the commuting-square tool, the obstruction-stratification tool, and the carrier-compression tool.
\end{enumerate}
\end{definition}

\begin{proposition}[Channel-sorted citation rule for the immediate continuation line]
\label{prop:channel-sorted-citation-rule}
Let a later argument remain inside the currently admitted OP1$\to$branch$\to$OP3/QAM$\to$OP5 architecture. Then the default citation discipline is as follows.
\begin{enumerate}[leftmargin=2em]
  \item If the argument only needs preservation of the retained total-channel profile or the subordination of visible diagnostics to one already fixed carrier package, it should cite the total-channel packet \(\mathfrak T_{\mathrm{tot}}\).
  \item If the argument accepts the retained total-channel surface as fixed and only compresses the surviving OP~1 burden toward the optimum-pinning kernel, it should cite the imbalance packet \(\mathfrak T_{\mathrm{imb}}\).
  \item If the argument must translate between total and imbalance language, or between channel language and a shell/menu/matrix realization, it should cite the bridge packet \(\mathfrak T_{\mathrm{br}}\).
\end{enumerate}
In particular, the default rule of the next revision is: reopen only the smallest packet that still sees the live burden.
\end{proposition}

\begin{proof}
This is exactly the burden split already extracted in the frozen line. The present proposition merely normalizes that split into a reusable citation discipline: total-channel stability is carried by the total packet, OP~1 compression by the imbalance packet, and all language-changing or layer-changing moves by the bridge packet.
\end{proof}

\begin{remark}[What the normalized ledger already gains]
\label{rem:normalized-ledger-gain}
The gain is immediate even before any new theorem is proved. First, repeated later arguments no longer need to cite a mixed list of total, imbalance, and translation tools in arbitrary order. Second, the surviving live burden is seen more honestly: OP~1 is not ``everything left,'' but specifically the active content of \(\mathfrak T_{\mathrm{imb}}\) once \(\mathfrak T_{\mathrm{tot}}\) has already been frozen. Third, the bridge packet isolates exactly where QAM, matrix, and cryptanalytic language should re-enter: only when one genuinely needs a change of surface.
\end{remark}

\subsection*{19AZ$'''$. QAM recoding of the normalized tool ledger}
\label{sec:qam-recoding-normalized-tool-ledger}
Once the channel-sorted ledger has been fixed, the right immediate QAM task is no longer to widen the public syntax blindly. It is to recode the already stabilized tool packets as a finite theorem-store/menu layer so that later proofs can query, combine, and route them without reopening the full historical derivation. This keeps QAM in its intended role: a definitional overlay and machine-readable proof interface, not a second independent mathematics.

\begin{definition}[QAM-coded tool menu]
\label{def:qam-coded-tool-menu}
Let \(\mathfrak T_{\mathrm{tot}}\), \(\mathfrak T_{\mathrm{imb}}\), and \(\mathfrak T_{\mathrm{br}}\) be the normalized packets from Definition~\ref{def:normalized-channel-sorted-tool-packets}. Their immediate QAM recoding is the finite coded menu
\[
\mathcal Q_{\mathrm{tool}}
:=
\bigl\{\langle \mathrm{tot},\tau\rangle : \tau\in\mathfrak T_{\mathrm{tot}}\bigr\}
\cup
\bigl\{\langle \mathrm{imb},\tau\rangle : \tau\in\mathfrak T_{\mathrm{imb}}\bigr\}
\cup
\bigl\{\langle \mathrm{br},\tau\rangle : \tau\in\mathfrak T_{\mathrm{br}}\bigr\}.
\]
Here the first coordinate records the burden class and the second coordinate records the already extracted uncoded tool. The purpose of \(\mathcal Q_{\mathrm{tool}}\) is indexing, routing, and proof compression; it does not create a new semantic layer beyond the underlying tool statements.
\end{definition}

\begin{definition}[Menu-to-carrier reduction]
\label{def:menu-to-carrier-reduction}
A current local proof obligation is said to admit a \emph{menu-to-carrier reduction} if, after coding the obligation into the finite menu \(\mathcal Q_{\mathrm{tool}}\), there exists a finite submenu \(M\subseteq \mathcal Q_{\mathrm{tool}}\) such that the underlying uncoded tools already imply either
\begin{enumerate}[leftmargin=2em]
  \item one bundled carrier-stable pass statement, or
  \item one strictly smaller localized obstruction class than the original uncoded presentation displayed.
\end{enumerate}
Thus the menu does not replace proof content; it compresses the route to the smallest carrier-bearing or obstruction-bearing interface.
\end{definition}

\begin{proposition}[QAM reuse principle for the immediate continuation line]
\label{prop:qam-reuse-principle-immediate-line}
If a current proof obligation admits menu-to-carrier reduction in the sense of Definition~\ref{def:menu-to-carrier-reduction}, then the next continuation line may process that obligation through the finite coded menu \(\mathcal Q_{\mathrm{tool}}\) instead of reopening the full local historical derivation. In particular, QAM now functions as the normalized theorem-store and routing overlay for the post-freeze tool layer.
\end{proposition}

\begin{proof}
By construction, every element of \(\mathcal Q_{\mathrm{tool}}\) is only a coded wrapper around an already extracted uncoded tool. If a finite submenu already implies the same bundled carrier statement or a strictly smaller obstruction class, then reopening the full derivation adds no new structural information. The coded menu therefore acts as a conservative proof-compression interface and not as a source of new independent mathematical content.
\end{proof}

\begin{remark}[Why this QAM step is mathematically honest]
\label{rem:qam-step-mathematically-honest}
This is the right scale for the immediate QAM revision. The present line does \emph{not} claim that the finite coded menu solves OP~1 or OP~5 automatically. It claims something narrower and more useful: once the main burden has already been split into total, imbalance, and bridge packets, QAM can store that split explicitly, route later obligations to the smallest relevant packet, and expose where proof compression is real rather than cosmetic.
\end{remark}


\subsection*{19AZ$''''$. First finite submenus and certified reuse routes inside the QAM tool menu}
\label{sec:first-finite-submenus-certified-reuse-routes}
The next optimization step is therefore to stop treating \(\mathcal Q_{\mathrm{tool}}\) as one undifferentiated finite store. For later proof work one wants a small number of \emph{canonical finite submenus} that already match the most common post-freeze obligations. These submenus do not add new mathematics; they only isolate the smallest coded interfaces that repeatedly suffice.

\begin{definition}[First canonical finite submenus of \(\mathcal Q_{\mathrm{tool}}\)]
\label{def:first-canonical-finite-submenus-qtool}
Inside the coded menu \(\mathcal Q_{\mathrm{tool}}\), define the following four canonical finite submenus.
\begin{enumerate}[leftmargin=2em]
  \item \textbf{Total-rigidity submenu} \(M^{\mathrm{rig}}_{\mathrm{tot}}\). This is the coded submenu generated by the transport-first tool, the readout-after-transport tool, the admissible-readout uniqueness tool, and the global pass/transfer/diagnostic-subordination interface already collected in the total-channel packet \(\mathfrak T_{\mathrm{tot}}\).

  \item \textbf{Optimum-pinning submenu} \(M^{\mathrm{pin}}_{\mathrm{imb}}\). This is the coded submenu generated by the frame-transport split certificate, the finite four-shell certificate, the two-pass half-burden certificate, the reciprocal pass-symmetry certificate, the same-corridor reversal implication, the unit-loop half-burden certificate, the one-first-zero loop-budget bridge, the admissible-window uniqueness bridge, the retained-window-rigidity / no-second-carrier / strict-upper-gap bridge, and the optimum-pinning comparison-rigidity bridge already collected in the imbalance packet \(\mathfrak T_{\mathrm{imb}}\).

  \item \textbf{Surface-change submenu} \(M^{\mathrm{surf}}_{\mathrm{br}}\). This is the coded submenu generated by the shared matrix-leakage tool, the cryptanalytic translation tool, the shared phase-coupling tool, the shell-stable readout tool, and the commuting-square tool from the bridge packet \(\mathfrak T_{\mathrm{br}}\). Its purpose is to move between channel language and shell/menu/matrix language without reopening the whole local history.

  \item \textbf{Obstruction-localization submenu} \(M^{\mathrm{loc}}_{\mathrm{br}}\). This is the coded submenu generated by the shared matrix-leakage tool, the shared phase-coupling tool, the shell-stable readout tool, the obstruction-stratification tool, and the carrier-compression tool from the same bridge packet. Its purpose is to decide where a still-live residual should be reopened first once the total-channel baseline is already frozen.
\end{enumerate}
\end{definition}

\begin{proposition}[First certified submenu reduction routes]
\label{prop:first-certified-submenu-reduction-routes}
Each of the following recurring post-freeze proof obligations admits menu-to-carrier reduction through one of the canonical finite submenus from Definition~\ref{def:first-canonical-finite-submenus-qtool}.
\begin{enumerate}[leftmargin=2em]
  \item Any obligation whose only job is to preserve the retained total-channel profile, carrier-stable pass status, or diagnostic subordination may be routed through \(M^{\mathrm{rig}}_{\mathrm{tot}}\).

  \item Any obligation that accepts the retained total-channel profile as fixed and only seeks to compress the surviving OP~1 burden toward retained-window rigidity or optimum pinning may be routed through \(M^{\mathrm{pin}}_{\mathrm{imb}}\).

  \item Any obligation whose main difficulty is a controlled change of surface --- from channel language to shell, menu, or matrix language, or back again --- may be routed through \(M^{\mathrm{surf}}_{\mathrm{br}}\).

  \item Any obligation whose main difficulty is not proof of a new carrier statement but localization of the first still-live residual interface may be routed through \(M^{\mathrm{loc}}_{\mathrm{br}}\).
\end{enumerate}
In particular, the immediate continuation line now has four explicit reuse routes rather than only one undifferentiated theorem store.
\end{proposition}

\begin{proof}
The total-rigidity submenu is by definition exactly the finite coded portion of \(\mathfrak T_{\mathrm{tot}}\) that already controls transport-first stability, readout-after-transport stability, admissible readout uniqueness, and global diagnostic subordination. Therefore every obligation of type \emph{(1)} is already reduced to the smallest currently relevant carrier-stable total-channel interface.

The optimum-pinning submenu is by definition the finite coded chain inside \(\mathfrak T_{\mathrm{imb}}\) that successively compresses the surviving OP~1 burden from frame split to shell balance, two-pass symmetry, loop budget, admissible-window uniqueness, retained-window rigidity, and finally optimum-pinning comparison-rigidity. Therefore every obligation of type \emph{(2)} is already reduced to the smallest currently relevant imbalance-side compression route.

The surface-change submenu isolates exactly those bridge tools whose role is transport across surfaces rather than new burden creation: matrix leakage, cryptanalytic translation, shared phase coupling, shell-stable readout, and commuting-square compatibility. Hence every obligation of type \emph{(3)} is reduced to the smallest currently relevant change-of-surface route.

Finally, the obstruction-localization submenu isolates the bridge tools whose role is not to certify a new global pass statement but to decide where the first still-live residual now sits: shared leakage/phase diagnostics identify the candidate bridge side, shell-stable readout preserves the common carrier-bearing content where available, obstruction stratification localizes the first failure stratum, and carrier compression collapses the readout/kernel side to one smallest missing carrier package. Hence every obligation of type \emph{(4)} admits menu-to-carrier reduction through \(M^{\mathrm{loc}}_{\mathrm{br}}\). This is exactly the claimed four-route reuse structure.
\end{proof}

\begin{corollary}[Working citation rule for the next immediate revision]
\label{cor:working-citation-rule-next-immediate-revision}
In the next immediate revision, a later proof should normally cite one of the four canonical finite submenus first and only reopen the larger packet description if that submenu fails to carry the obligation. Concretely:
\begin{enumerate}[leftmargin=2em]
  \item cite \(M^{\mathrm{rig}}_{\mathrm{tot}}\) first for total-channel rigidity or status-preservation steps;
  \item cite \(M^{\mathrm{pin}}_{\mathrm{imb}}\) first for OP~1 compression toward the optimum-pinning kernel;
  \item cite \(M^{\mathrm{surf}}_{\mathrm{br}}\) first for channel/shell/menu/matrix translation steps;
  \item cite \(M^{\mathrm{loc}}_{\mathrm{br}}\) first for localization of the first still-live bridge or readout-side obstruction.
\end{enumerate}
Thus the QAM layer has now moved from abstract finite storage to a concrete four-route routing discipline.
\end{corollary}

\begin{remark}[What has been gained before proving anything new]
\label{rem:what-gained-before-proving-anything-new-submenus}
The gain is practical and immediate. A later argument no longer has to say only ``use QAM'' or ``reopen the bridge packet.'' It can now say much more sharply: preserve the frozen total profile via \(M^{\mathrm{rig}}_{\mathrm{tot}}\), compress the surviving OP~1 route via \(M^{\mathrm{pin}}_{\mathrm{imb}}\), switch surfaces via \(M^{\mathrm{surf}}_{\mathrm{br}}\), or localize the first still-live residual via \(M^{\mathrm{loc}}_{\mathrm{br}}\). That is exactly the kind of proof-shortening interface the immediate tool/QAM revision was supposed to extract.
\end{remark}

\subsection*{19AZ$^{(3a)}$. Normalized QAM nomenclature and finite tag layer for the four canonical submenus}
\label{sec:normalized-qam-tag-layer}
The four canonical submenus should now receive a short address language. Without that extra layer, the current QAM revision would still have finite storage and routing, but not yet a stable nomenclature by which later proofs or diagnostics can point to one smallest already-normalized carrier without reopening the full sentence around it.

\begin{definition}[Primary tool-tag families]
\label{def:primary-tool-tag-families}
Read the four canonical finite submenus of Definition~\ref{def:qam-coded-tool-menu} in the order induced by their listed entries. The \emph{primary tool-tag families} are the finite code families
\[
\mathrm{TR}_1,\mathrm{TR}_2,\dots;\qquad
\mathrm{IP}_1,\mathrm{IP}_2,\dots;\qquad
\mathrm{BS}_1,\mathrm{BS}_2,\dots;\qquad
\mathrm{BL}_1,\mathrm{BL}_2,\dots
\]
indexed respectively by the ordered entries of
\[
M^{\mathrm{rig}}_{\mathrm{tot}},\qquad
M^{\mathrm{pin}}_{\mathrm{imb}},\qquad
M^{\mathrm{surf}}_{\mathrm{br}},\qquad
M^{\mathrm{loc}}_{\mathrm{br}}.
\]
Here \(\mathrm{TR}\) stands for \emph{total rigidity}, \(\mathrm{IP}\) for \emph{imbalance pinning}, \(\mathrm{BS}\) for \emph{bridge/surface change}, and \(\mathrm{BL}\) for \emph{bridge/localization}. The ordered disjoint union of these four families is the \emph{primary tool-tag alphabet} of the immediate continuation line.
\end{definition}

\begin{definition}[Primary tool-tag registry and dereferencing map]
\label{def:primary-tool-tag-registry-dereferencing-map}
The \emph{primary tool-tag registry}
\[
\mathcal R^{\mathrm{tag}}_{\mathrm{tool}}
\]
is the finite family of quadruples
\[
(\tau,M,\kappa,x)
\]
with the following meaning:
\begin{enumerate}[leftmargin=2em]
  \item \(\tau\) is one primary tag from Definition~\ref{def:primary-tool-tag-families};
  \item \(M\) is the unique canonical submenu among \(M^{\mathrm{rig}}_{\mathrm{tot}}, M^{\mathrm{pin}}_{\mathrm{imb}}, M^{\mathrm{surf}}_{\mathrm{br}}, M^{\mathrm{loc}}_{\mathrm{br}}\) in which the tagged item has its primary home;
  \item \(\kappa\in\{\mathrm{status},\mathrm{pinning},\mathrm{surface},\mathrm{localization}\}\) is its normalized burden type;
  \item \(x\in M\) is the unique carrier statement, submenu entry, or smallest normalized reusable route addressed by \(\tau\).
\end{enumerate}
The induced map
\[
\mathsf{Ref}:\tau\longmapsto x
\]
is called the \emph{dereferencing map}. If one and the same carrier statement is reusable from more than one submenu, then exactly one of those occurrences is designated as its \emph{primary home}; every other appearance is recorded only as a secondary pointer and receives no second primary tag. The registry order defines a precedence rank
\[
\pi(\tau)\in\mathbb N
\]
inside the chosen submenu, so that smaller rank means earlier reopening inside the already chosen route.
\end{definition}

\begin{proposition}[Tag soundness, uniqueness, and minimal reopening discipline]
\label{prop:tag-soundness-uniqueness-minimal-reopening-discipline}
For every entry \(x\) of the four canonical finite submenus there exists exactly one primary tag \(\tau\) such that
\[
(\tau,M,\kappa,x)\in \mathcal R^{\mathrm{tag}}_{\mathrm{tool}}.
\]
Consequently:
\begin{enumerate}[leftmargin=2em]
  \item the primary tag \(\tau\) determines the canonical submenu \(M\), the normalized burden type \(\kappa\), and the dereferenced carrier \(x\);
  \item no collision between the four nomenclature families can occur at the primary level;
  \item once a later argument has chosen its submenu, reopening should proceed in increasing precedence rank \(\pi(\tau)\) until the first sufficient carrier is reached.
\end{enumerate}
Hence the immediate QAM layer now carries a finite nomenclature with exact dereferencing and a smallest-first reopening discipline.
\end{proposition}

\begin{proof}
The four canonical submenus are finite and ordered by construction. Assigning one tag from the corresponding family to each ordered entry therefore produces exactly one registry quadruple \((\tau,M,\kappa,x)\) for each such entry. Because the tag families \(\{\mathrm{TR}_i\}, \{\mathrm{IP}_j\}, \{\mathrm{BS}_k\}, \{\mathrm{BL}_\ell\}\) are disjoint by notation, primary collisions cannot occur across different submenu types. If one carrier statement appears in more than one submenu role, the definition of primary home forces uniqueness at the registry level by declaring every other occurrence secondary only. Finally, the submenu order itself supplies the precedence rank \(\pi(\tau)\), so once a route has been chosen the smallest-first reopening order is already fixed. This is exactly the claimed tag soundness, uniqueness, and reopening discipline.
\end{proof}

\begin{corollary}[Tag-first citation syntax for the immediate continuation line]
\label{cor:tag-first-citation-syntax-immediate-continuation-line}
The next immediate revision may cite the normalized tool layer in either of the following two equivalent ways:
\begin{enumerate}[leftmargin=2em]
  \item by naming the submenu and then the primary tag, for example ``invoke \(M^{\mathrm{rig}}_{\mathrm{tot}}\) at \(\mathrm{TR}_i\)'' or ``reopen \(M^{\mathrm{loc}}_{\mathrm{br}}\) at \(\mathrm{BL}_\ell\)'';
  \item once the submenu is fixed by context, by citing the primary tag alone and dereferencing it through \(\mathsf{Ref}\).
\end{enumerate}
Accordingly, the four primary tag families should be read as the first normalized QAM nomenclature layer for the immediate continuation line:
\begin{center}
\begin{tabular}{@{}ll@{}}
\toprule
Tag family & Intended immediate use \\
\midrule
\(\mathrm{TR}_i\) & total-channel rigidity and status preservation \\
\(\mathrm{IP}_j\) & OP~1 compression toward the optimum-pinning kernel \\
\(\mathrm{BS}_k\) & channel/shell/menu/matrix surface change \\
\(\mathrm{BL}_\ell\) & localization of the first still-live obstruction \\
\bottomrule
\end{tabular}
\end{center}
Thus the QAM layer now has not only finite menus but also a short stable address language for them.
\end{corollary}

\begin{remark}[What the tag layer buys before any new theorem is added]
\label{rem:what-tag-layer-buys-before-new-theorem}
This is the first genuinely addressable QAM tool surface on the present line. A later proof can now point to a smallest normalized carrier by a short typed tag instead of reopening a long theorem sentence; a later diagnostic routine can record that a residual cluster repeatedly points toward \(\mathrm{IP}_7\) or \(\mathrm{BL}_2\) rather than toward an entire section; and a later public extraction can export the tags as the front-face nomenclature without exporting the whole internal packet architecture. In that precise sense, the tag layer is the missing bridge from finite storage to reusable proof vocabulary.
\end{remark}



\subsection*{19AZ$^{(3b)}$. First explicit primary registry assignments for the normalized tool tags}
\label{sec:first-explicit-primary-registry-assignments}
The normalized tag alphabet should now be instantiated once, so that the immediate continuation line can cite concrete dereferenced carriers instead of only a generic family pattern. At the present stage the assignments are intentionally conservative: each tag points only to a carrier statement or smallest reusable route that already exists in the frozen working text.

\paragraph{First concrete primary registry.}
Write the finite registry in four family blocks. For compactness, the ``carrier'' column records the dereferenced item together with its smallest currently stable reference surface.
\begin{quote}\small
\noindent\textbf{Reader cue.} Read the following registry tables as lookup surfaces for the normalized QAM tags. Their purpose is to pin short addresses to already stabilized statements or reusable micro-routes; they do not introduce new carriers.
\end{quote}

{\small
\renewcommand{\arraystretch}{1.08}
\begin{center}
\begin{tabular}{@{}>{\raggedright\arraybackslash}p{0.10\linewidth}>{\raggedright\arraybackslash}p{0.14\linewidth}>{\raggedright\arraybackslash}p{0.66\linewidth}@{}}
\toprule
Tag & $\kappa$ & Dereferenced carrier \\
\midrule
\multicolumn{3}{@{}l@{}}{\textbf{Total-rigidity submenu} $M^{\mathrm{rig}}_{\mathrm{tot}}$}\\[0.2em]
$\mathrm{TR}_1$ & status & the transport-first route read from the opening clause of the proof of Theorem~\ref{thm:working-op1-admissible-readout-uniqueness}, namely: branch selection is fixed before visible readout is interpreted;\\
$\mathrm{TR}_2$ & status & the readout-after-transport route read from the next clause of the same proof of Theorem~\ref{thm:working-op1-admissible-readout-uniqueness}, namely: every canonical readout remains downstream from the already fixed transport branch;\\
$\mathrm{TR}_3$ & status & the admissible-readout uniqueness carrier of Theorem~\ref{thm:working-op1-admissible-readout-uniqueness};\\
$\mathrm{TR}_4$ & status & the bundled global pass/transfer/diagnostic-subordination interface from Propositions~\ref{prop:global-pass-tool}, \ref{prop:global-transfer-tool}, and \ref{prop:global-diagnostic-subordination-tool}.\\
\bottomrule
\end{tabular}
\end{center}

\begin{center}
\begin{tabular}{@{}>{\raggedright\arraybackslash}p{0.10\linewidth}>{\raggedright\arraybackslash}p{0.14\linewidth}>{\raggedright\arraybackslash}p{0.66\linewidth}@{}}
\toprule
Tag & $\kappa$ & Dereferenced carrier \\
\midrule
\multicolumn{3}{@{}l@{}}{\textbf{Optimum-pinning submenu} $M^{\mathrm{pin}}_{\mathrm{imb}}$}\\[0.2em]
$\mathrm{IP}_1$ & pinning & frame-transport split certificate, Proposition~\ref{prop:op1-frame-transport-split};\\
$\mathrm{IP}_2$ & pinning & finite four-shell certificate, Proposition~\ref{prop:op1-four-shell-certificate};\\
$\mathrm{IP}_3$ & pinning & two-pass half-burden certificate, Proposition~\ref{prop:op1-two-pass-half-burden};\\
$\mathrm{IP}_4$ & pinning & reciprocal pass-symmetry certificate, Proposition~\ref{prop:op1-reciprocal-pass-symmetry};\\
$\mathrm{IP}_5$ & pinning & same-corridor reversal implication, Proposition~\ref{prop:op1-same-corridor-implies-reciprocity};\\
$\mathrm{IP}_6$ & pinning & unit-loop half-burden certificate, Proposition~\ref{prop:op1-unit-loop-half-burden};\\
$\mathrm{IP}_7$ & pinning & one-first-zero loop-budget bridge, Proposition~\ref{prop:op1-one-first-zero-loop-budget};\\
$\mathrm{IP}_8$ & pinning & admissible-window uniqueness bridge, Proposition~\ref{prop:op1-admissible-window-uniqueness};\\
$\mathrm{IP}_9$ & pinning & retained-window-rigidity / no-second-carrier / strict-upper-gap bridge, Proposition~\ref{prop:op1-rigidity-gap-uniqueness-bridge};\\
$\mathrm{IP}_{10}$ & pinning & optimum-pinning comparison-rigidity bridge, Proposition~\ref{prop:op1-optimum-pinning-implies-rigidity}.\\
\bottomrule
\end{tabular}
\end{center}

\begin{center}
\begin{tabular}{@{}>{\raggedright\arraybackslash}p{0.10\linewidth}>{\raggedright\arraybackslash}p{0.14\linewidth}>{\raggedright\arraybackslash}p{0.66\linewidth}@{}}
\toprule
Tag & $\kappa$ & Dereferenced carrier \\
\midrule
\multicolumn{3}{@{}l@{}}{\textbf{Surface-change submenu} $M^{\mathrm{surf}}_{\mathrm{br}}$}\\[0.2em]
$\mathrm{BS}_1$ & surface & shared matrix-leakage surface gate, read through Definition~\ref{def:working-shared-matrix-leakage} together with Proposition~\ref{prop:working-stable-channel-criterion};\\
$\mathrm{BS}_2$ & surface & the cryptanalytic translation route from the shared matrix-leakage / cryptanalytic block;\\
$\mathrm{BS}_3$ & surface & shared phase-coupling surface diagnostic from Definitions~\ref{def:working-shared-phase-traces}--\ref{def:working-common-carrier-diagnostic} together with Propositions~\ref{prop:working-one-mode-image-test} and \ref{prop:working-quadrature-criterion};\\
$\mathrm{BS}_4$ & surface & the shell-stable readout tool from the extracted OP~3/QAM tool layer;\\
$\mathrm{BS}_5$ & surface & the commuting-square tool from the extracted OP~3/QAM tool layer.\\
\bottomrule
\end{tabular}
\end{center}

\begin{center}
\begin{tabular}{@{}>{\raggedright\arraybackslash}p{0.10\linewidth}>{\raggedright\arraybackslash}p{0.14\linewidth}>{\raggedright\arraybackslash}p{0.66\linewidth}@{}}
\toprule
Tag & $\kappa$ & Dereferenced carrier \\
\midrule
\multicolumn{3}{@{}l@{}}{\textbf{Obstruction-localization submenu} $M^{\mathrm{loc}}_{\mathrm{br}}$}\\[0.2em]
$\mathrm{BL}_1$ & localization & matrix-leakage-first localization gate toward the underfixed HC/MM/J coupling, read again through Definition~\ref{def:working-shared-matrix-leakage} and Proposition~\ref{prop:working-stable-channel-criterion};\\
$\mathrm{BL}_2$ & localization & phase-coupling-first localization gate asking whether the residual pair is still one masked carrier, read through Propositions~\ref{prop:working-one-mode-image-test} and \ref{prop:working-quadrature-criterion};\\
$\mathrm{BL}_3$ & localization & the shell-stable readout localization gate from the extracted OP~3/QAM tool layer;\\
$\mathrm{BL}_4$ & localization & the obstruction-stratification tool from the extracted OP~3/QAM tool layer;\\
$\mathrm{BL}_5$ & localization & the carrier-compression tool from the extracted OP~3/QAM tool layer.\\
\bottomrule
\end{tabular}
\end{center}
}

\paragraph{Primary-home convention at the concrete level.}
The bridge families contain repeated antecedent references on purpose. This does not violate the primary-home rule from Definition~\ref{def:primary-tool-tag-registry-dereferencing-map}, because the registry items above are not identified merely by their antecedent theorem labels, but by their \emph{normalized route role}. Thus the same matrix-leakage or shell-stable-readout material may appear once as a \emph{surface-change gate} and once as a \emph{localization gate} without creating a primary collision.

\paragraph{Immediate working use.}
From this point onward the continuation line may cite concrete tags without any placeholder notation. For example, ``freeze the total channel at $\mathrm{TR}_4$'', ``compress the surviving OP~1 kernel through $\mathrm{IP}_8\to \mathrm{IP}_{10}$'', ``change surface via $\mathrm{BS}_2$'', or ``localize the first still-live residual by testing $\mathrm{BL}_1$, then $\mathrm{BL}_2$, then $\mathrm{BL}_4$'' are now literal dereferenceable instructions on the QAM side.


\subsection*{19AZ$^{(3c)}$. First tag-first rewrite of the surviving OP~1 route}
\label{sec:first-tag-first-rewrite-of-the-surviving-op1-route}
The normalized registry should now be used once on a real proof route rather than remaining only a citation promise. The cleanest first test is the currently surviving OP~1 compression chain, because it already has a sharply localized terminal burden and because its hypotheses were precisely the ones that motivated the split into total, pinning, and bridge roles.

\begin{proposition}[First tag-first reformulation of the surviving OP~1 closure route]
\label{prop:first-tag-first-op1-route}
Assume the same-corridor reversal certificate of Definition~\ref{def:op1-same-corridor-reversal}, and assume moreover that the optimum-pinning comparison-rigidity certificate of Definition~\ref{def:op1-optimum-pinning-comparison-rigidity} holds on the intended admissible windows. Then the surviving OP~1 closure route may be cited on the normalized QAM side in the compact form
\[
  \mathrm{TR}_4 \;\Longrightarrow\; \mathrm{IP}_{10} \;\Longrightarrow\; \mathrm{IP}_9 \;\Longrightarrow\; \mathrm{IP}_8,
\]
and therefore
\[
  \Delta_{\mathrm{CL}} = 0.
\]
Equivalently, once the retained total-channel baseline and its bundled transfer/diagnostic discipline are frozen through \(\mathrm{TR}_4\), the live OP~1 burden is legitimately reopened only as the short pinning chain
\[
  \mathrm{IP}_{10} \to \mathrm{IP}_9 \to \mathrm{IP}_8,
\]
rather than as a return to the full historical OP~1 queue.
\end{proposition}

\begin{proof}
By the concrete registry from Section~\ref{sec:first-explicit-primary-registry-assignments}, the tag \(\mathrm{TR}_4\) dereferences the bundled global pass/transfer/diagnostic-subordination interface from Propositions~\ref{prop:global-pass-tool}, \ref{prop:global-transfer-tool}, and \ref{prop:global-diagnostic-subordination-tool}. On the present line this is exactly the smallest normalized total-channel surface that keeps fixed the retained total profile and the comparison/diagnostic discipline under which the OP~1 corridor is being read.

Under the standing same-corridor reversal hypothesis, the additional assumption of optimum-pinning comparison rigidity is precisely the carrier tagged by \(\mathrm{IP}_{10}\), namely Proposition~\ref{prop:op1-optimum-pinning-implies-rigidity}. That proposition says that, on the intended admissible windows, preservation of the retained total-channel profile together with the retained optimum location and local comparison/defect data forces retained-window rigidity.

The next normalized reopening step is then \(\mathrm{IP}_9\), i.e. Proposition~\ref{prop:op1-rigidity-gap-uniqueness-bridge}. Once retained-window rigidity is available, and with the standing no-second-carrier and strict-upper-gap exclusions built into that bridge, admissible-window uniqueness follows on the retained corridor.

Finally one reopens \(\mathrm{IP}_8\), i.e. Proposition~\ref{prop:op1-admissible-window-uniqueness}. This upgrades admissible-window uniqueness to the one-first-zero loop-budget certificate and hence yields
\[
  \Delta_{\mathrm{CL}} = 0.
\]
No new mathematics has been added here: the content is exactly the already established OP~1 compression chain, now rewritten in normalized tag-first form. The gain is that the route can be cited as one frozen total-channel anchor plus three smallest pinning tags instead of reopening the larger sentence-level proof history.
\end{proof}

\begin{remark}[What this first practical rewrite already buys]
\label{rem:what-the-first-practical-tag-rewrite-buys}
This is the first point at which the tag layer stops being merely architectural. A later continuation can now say, in a mathematically literal way, ``freeze at \(\mathrm{TR}_4\), then test \(\mathrm{IP}_{10}\to\mathrm{IP}_9\to\mathrm{IP}_8\)'' and thereby refer to an already normalized proof route rather than to an entire subsection block. In other words, the QAM layer has started to function as a true routing vocabulary for the live burden, not just as a storage format for tools.
\end{remark}


\subsection*{19AZ$^{(3d)}$. First tag-first rewrite of the bridge/localization route}
\label{sec:first-tag-first-bridge-localization-route}
The second practical test of the normalized QAM registry should not pretend to close a new theorematic burden. Its job is narrower and more honest: once the retained total-channel baseline is frozen, it should rewrite the currently live bridge-side reopening order into one compact route that first changes surface at the smallest matrix/phase interface and only then localizes the first still-live obstruction.

\begin{proposition}[First tag-first reformulation of the bridge/localization reopening route]
\label{prop:first-tag-first-bridge-localization-route}
Assume the retained total-channel baseline is frozen through \(\mathrm{TR}_4\), and suppose a still-live bridge-side residual is to be tested without yet postulating two independent primitive carriers. Then the first normalized QAM reopening route may be cited in the compact form
\[
  \mathrm{TR}_4 \;\Longrightarrow\; \mathrm{BS}_1 \;\Longrightarrow\; \mathrm{BS}_3 \;\Longrightarrow\; \mathrm{BL}_2 \;\Longrightarrow\; \mathrm{BL}_4,
\]
with the additional compression step
\[
  \mathrm{BL}_4 \;\Longrightarrow\; \mathrm{BL}_5
\]
whenever the first localized obstruction lies in the readout/kernel stratum.
Equivalently, once the retained total profile and its bundled comparison/diagnostic discipline are frozen, the live bridge burden is reopened first as matrix leakage, then as common-carrier/phase coupling, and only then as stratified obstruction localization; a separate carrier is not introduced before this normalized route has been exhausted.
\end{proposition}

\begin{proof}
By the concrete registry from Section~\ref{sec:first-explicit-primary-registry-assignments}, the tag \(\mathrm{TR}_4\) dereferences the bundled global pass/transfer/diagnostic-subordination interface from Propositions~\ref{prop:global-pass-tool}, \ref{prop:global-transfer-tool}, and \ref{prop:global-diagnostic-subordination-tool}. On the present line this is again the smallest frozen total-channel surface from which later local re-openings are allowed.

The first bridge-side reopening step is \(\mathrm{BS}_1\), i.e. the shared matrix-leakage surface gate from Definition~\ref{def:working-shared-matrix-leakage} together with Proposition~\ref{prop:working-stable-channel-criterion}. This is the right first test because it asks whether there is still matrix-induced off-sector drift at all. In other words, before any finer localization language is used, the route first checks whether the admissible HC channel is already stable under the Millennium/Umkehrwalze interface.

If the bridge burden is still live after that first gate, then one changes surface through \(\mathrm{BS}_3\), namely the shared phase-coupling diagnostic from Definitions~\ref{def:working-shared-phase-traces}--\ref{def:working-common-carrier-diagnostic} together with Propositions~\ref{prop:working-one-mode-image-test} and \ref{prop:working-quadrature-criterion}. This does not yet decide where the first remaining obstruction sits; it rewrites the visible residual pair into the smallest matrix/phase language in which one can honestly ask whether the pair is still one masked carrier rather than two unrelated end defects.

That localization question is then reopened directly as \(\mathrm{BL}_2\), which is precisely the phase-coupling-first localization gate recorded in the same registry. If the residual pair still behaves as one masked carrier, the route stays on the one-carrier side. If it remains unresolved or fails at that stage, the next legitimate reopening is \(\mathrm{BL}_4\), i.e. the obstruction-stratification tool from the extracted OP~3/QAM layer. This is the first point at which the present line is allowed to ask whether the surviving obstruction lives at the branch, shell, or readout/kernel level.

Finally, when \(\mathrm{BL}_4\) places the first residual in the readout/kernel stratum, the carrier-compression tool \(\mathrm{BL}_5\) compresses that stratum to absence of one common shell-stable carrier. Hence the entire normalized reopening order is exactly the displayed tag route. No new closure theorem has been proved here. The content is the already extracted bridge packet and the already extracted localization packet, now rewritten as one smallest admissible tag-first diagnostic route.
\end{proof}

\begin{remark}[What this second practical rewrite already buys]
\label{rem:what-the-second-practical-tag-rewrite-buys}
The present line can now cite both kinds of live work in the same compact language. For OP~1 compression one may say ``freeze at \(\mathrm{TR}_4\), then test \(\mathrm{IP}_{10}\to\mathrm{IP}_9\to\mathrm{IP}_8\)''; for bridge/localization one may now say ``freeze at \(\mathrm{TR}_4\), then test \(\mathrm{BS}_1\to\mathrm{BS}_3\), then localize by \(\mathrm{BL}_2\to\mathrm{BL}_4\), and if necessary compress by \(\mathrm{BL}_5\).'' In that precise sense, the QAM layer is no longer only a menu of stored tools: it has become a two-sided routing vocabulary, one side for the surviving OP~1 pinning chain and one side for the still-live bridge/localization diagnostics.
\end{remark}

\subsection*{19AZ$^{(4)}$. First Monte-Carlo residual-signature route beneath the normalized QAM tool packets}
\label{sec:first-monte-carlo-residual-signature-route}
The Monte-Carlo programme should therefore enter only \emph{after} the normalized ledger and its QAM recoding are in place. On the present line, Monte Carlo is not a replacement proof engine. It is a structural probe that samples how residual deviations distribute relative to the already fixed total-channel baseline and to the still-live imbalance/bridge packets.

\begin{definition}[Residual-signature word]
\label{def:residual-signature-word}
Fix a retained admissible window and a finite family of normalized diagnostics
\[
D=(D_1,\dots,D_m)
\]
measured relative to the retained total-channel baseline. For one Monte-Carlo sample \(\omega\), let
\[
\delta(\omega)=(\delta_1(\omega),\dots,\delta_m(\omega))
\]
be the resulting normalized deviation vector. A \emph{residual-signature word} for \(\omega\) is any finite coded word
\[
\mathsf{Sig}(\omega;D)
\]
obtained from the ordered sign, bin, rank, or slot pattern of the components of \(\delta(\omega)\). Such a word is interpreted as a structural signature or hash of the sampled residual behavior, not as a theorem by itself.
\end{definition}

\begin{definition}[Kasiski-style recurrence scan]
\label{def:kasiski-style-recurrence-scan}
Given a family of residual-signature words, the \emph{Kasiski-style recurrence scan} is the first combinatorial pass that looks for repeated blocks, repeated slot offsets, and stable periodicity or near-periodicity classes among those words. Its role is diagnostic only: it asks whether the sampled residuals point repeatedly toward the same bridge or imbalance surface.
\end{definition}

\begin{proposition}[Conservative Monte-Carlo/QAM probe rule]
\label{prop:conservative-monte-carlo-qam-probe-rule}
Residual-signature words and their Kasiski-style recurrence scan may be used in the immediate continuation line only to prioritize which normalized packet from Definition~\ref{def:normalized-channel-sorted-tool-packets} should be reopened first. They may \emph{not} by themselves introduce a second primitive carrier, overturn the frozen total-channel baseline, or replace an already theorematic pass/obstruction statement.
\end{proposition}

\begin{proof}
This is exactly the methodological discipline stated in the deferred roadmap above. The retained total-channel baseline is already part of the frozen theorem layer, while the Monte-Carlo route is introduced only as a structural probe beneath that layer. Therefore sampled residual signatures may guide priority and localization, but they cannot overrule the already established carrier and status surfaces.
\end{proof}

\begin{remark}[Immediate diagnostic reading of the first residual signatures]
\label{rem:immediate-diagnostic-reading-signatures}
The intended first reading is now precise. If the residual signatures cluster while the retained total-channel profile stays rigid, then the next reopen candidate should be searched first inside \(\mathfrak T_{\mathrm{imb}}\) or \(\mathfrak T_{\mathrm{br}}\), depending on whether the recurrence is purely local/imbalance-side or genuinely translation-sensitive. If no stable recurrence class appears at all, then the theorem layer should remain untouched and the Monte-Carlo route should be treated as presently non-informative rather than overinterpreted.
\end{remark}


\subsection*{19AZ$^{(4a)}$. First decoding of residual-signature recurrences into normalized BS/BL reopening priorities}
\label{sec:first-decoding-of-signature-recurrences-into-bs-bl-priorities}
The next conservative step is now to make the bridge/localization reading operational. The residual-signature words should not merely say that ``something bridge-like remains.'' They should decode into a smallest prioritized reopening order. In particular, the present line needs one first rule saying when a stable recurrence class should push the next diagnostic pass onto the normalized bridge/localization route rather than back into the OP~1 pinning chain.

\begin{definition}[Route-priority decoder for residual signatures]
\label{def:route-priority-decoder-residual-signatures}
Let \(\mathcal C\) be a stable recurrence class produced by the Kasiski-style scan from Definition~\ref{def:kasiski-style-recurrence-scan}. A \emph{route-priority decoder} for \(\mathcal C\) is any conservative rule
\[
\mathsf{Dec}(\mathcal C)\in \bigl\{M^{\mathrm{pin}}_{\mathrm{imb}},\,M^{\mathrm{surf}}_{\mathrm{br}},\,M^{\mathrm{loc}}_{\mathrm{br}}\bigr\}
\]
that assigns \(\mathcal C\) to the smallest normalized menu whose recorded carriers already account for the repeated slot/rank pattern of the sampled residuals.
The decoder is called \emph{translation-sensitive} if the repeated pattern is not stable under purely local imbalance-side rereading, but stabilizes once the residual pair is reread through the matrix/phase common-carrier surface. In that case the first prioritized reopening route attached to \(\mathcal C\) is
\[
\mathrm{BS}_1 \to \mathrm{BS}_3 \to \mathrm{BL}_2,
\]
with \(\mathrm{BL}_4\) and, when needed, \(\mathrm{BL}_5\) appended only after the phase/common-carrier gate has been exhausted.
\end{definition}

\begin{proposition}[Translation-sensitive signature classes prioritize the BS/BL route]
\label{prop:translation-sensitive-signature-classes-prioritize-bs-bl-route}
Assume the retained total-channel baseline remains frozen through \(\mathrm{TR}_4\), and let \(\mathcal C\) be a stable recurrence class of residual-signature words. Suppose moreover that \(\mathcal C\) is translation-sensitive in the sense of Definition~\ref{def:route-priority-decoder-residual-signatures}: its repeated slot/rank pattern does not settle under a purely local imbalance-side rereading, but does settle once the same sampled residuals are reread through the matrix/phase common-carrier surface. Then the first normalized reopening priority attached to \(\mathcal C\) is the bridge/localization route
\[
\mathrm{TR}_4 \;\Longrightarrow\; \mathrm{BS}_1 \;\Longrightarrow\; \mathrm{BS}_3 \;\Longrightarrow\; \mathrm{BL}_2,
\]
followed, only if a still-live obstruction remains, by
\[
\mathrm{BL}_2 \;\Longrightarrow\; \mathrm{BL}_4,
\]
and by
\[
\mathrm{BL}_4 \;\Longrightarrow\; \mathrm{BL}_5
\]
whenever the first localized obstruction lies in the readout/kernel stratum.
Equivalently, a stable translation-sensitive recurrence class is not yet evidence for a second primitive carrier; it is evidence that the next reopen candidate should be searched first along the normalized matrix/phase-to-localization bridge route.
\end{proposition}

\begin{proof}
By Proposition~\ref{prop:conservative-monte-carlo-qam-probe-rule}, residual-signature words may only prioritize a reopening order inside the already normalized theorem/tool layer; they may not overrule the frozen total-channel baseline or create a new carrier hypothesis. Hence one must begin from \(\mathrm{TR}_4\), the smallest bundled total-channel anchor already fixed by the registry from Section~\ref{sec:first-explicit-primary-registry-assignments}.

Now assume the stable recurrence class \(\mathcal C\) is translation-sensitive. By definition this means that the repeated residual pattern does not become stable when reread only as a local imbalance perturbation, but does become stable once the same samples are reread through the matrix/phase common-carrier surface. Therefore the first admissible decoder target for \(\mathcal C\) cannot be merely \(M^{\mathrm{pin}}_{\mathrm{imb}}\); the smallest normalized reopening surface that can already account for the stable repeated pattern is the bridge-surface menu. In the concrete registry this is exactly the route starting at \(\mathrm{BS}_1\), the shared matrix-leakage gate, and continuing through \(\mathrm{BS}_3\), the shared phase/common-carrier surface.

Once that surface change has been made, the first honest question is whether the repeated residual pair still behaves as one masked carrier. In the normalized registry this is precisely the localization gate \(\mathrm{BL}_2\). If the recurrence class is already resolved there, no deeper reopening is justified. If it remains live, then the first further admissible step is \(\mathrm{BL}_4\), the obstruction-stratification tool. And if that stratification places the first obstruction in the readout/kernel stratum, then \(\mathrm{BL}_5\) is exactly the recorded compression step that rewrites the surviving burden as absence of one common shell-stable carrier package. Thus the displayed BS/BL priority route is exactly the smallest normalized decoding of a translation-sensitive stable recurrence class.
\end{proof}

\begin{remark}[What the first Monte-Carlo-to-route decoding now buys]
\label{rem:what-the-first-monte-carlo-to-route-decoding-buys}
The Monte-Carlo layer now has a literal output format inside the QAM vocabulary. A stable recurrence class is no longer only a vague hint that ``the bridge side may matter.'' The present line may now state, in compact normalized form: ``freeze at \(\mathrm{TR}_4\); if the recurrence is translation-sensitive, decode first by \(\mathrm{BS}_1\to\mathrm{BS}_3\to\mathrm{BL}_2\), and only then localize further by \(\mathrm{BL}_4\) or compress by \(\mathrm{BL}_5\).'' In that precise sense the residual-signature words have acquired a conservative but genuine routing semantics.
\end{remark}



\subsection*{19AZ$^{(4b)}$. First decoding of purely local residual-signature recurrences into normalized IP reopening priorities}
\label{sec:first-decoding-of-local-signature-recurrences-into-ip-priorities}
The decoder should now be made symmetric on the pinning side. The residual-signature words should not only tell us when to reopen the bridge/localization route; they should also tell us when the next honest move remains entirely inside the surviving OP~1 kernel. The immediate line therefore needs one first rule saying when a stable recurrence class should stay on the local imbalance side and decode directly into the normalized IP route.

\begin{definition}[Local-imbalance decoder for residual signatures]
\label{def:local-imbalance-decoder-residual-signatures}
Let \(\mathcal C\) be a stable recurrence class produced by the Kasiski-style scan from Definition~\ref{def:kasiski-style-recurrence-scan}. A route-priority decoder from Definition~\ref{def:route-priority-decoder-residual-signatures} is called \emph{purely local} or \emph{imbalance-stable} on \(\mathcal C\) if the repeated slot/rank pattern already stabilizes under purely local rereading on the imbalance side, so that rereading the same samples through the matrix/phase common-carrier surface adds no smaller admissible explanation. In that case the first prioritized reopening route attached to \(\mathcal C\) is the normalized pinning chain
\[
\mathrm{IP}_{10} \to \mathrm{IP}_9 \to \mathrm{IP}_8,
\]
read beneath the frozen total-channel anchor \(\mathrm{TR}_4\).
\end{definition}

\begin{proposition}[Purely local signature classes prioritize the IP route]
\label{prop:purely-local-signature-classes-prioritize-ip-route}
Assume the retained total-channel baseline remains frozen through \(\mathrm{TR}_4\), and let \(\mathcal C\) be a stable recurrence class of residual-signature words. Suppose moreover that \(\mathcal C\) is purely local in the sense of Definition~\ref{def:local-imbalance-decoder-residual-signatures}: its repeated slot/rank pattern already stabilizes under a local imbalance-side rereading and does not require an earlier matrix/phase surface change in order to become coherent. Then the first normalized reopening priority attached to \(\mathcal C\) is the pinning route
\[
\mathrm{TR}_4 \;\Longrightarrow\; \mathrm{IP}_{10} \;\Longrightarrow\; \mathrm{IP}_9 \;\Longrightarrow\; \mathrm{IP}_8,
\]
so that the sampled recurrence is read first as a compressed OP~1 burden rather than as a bridge/localization burden.
Equivalently, a stable purely local recurrence class is not evidence that one should leave the imbalance side; it is evidence that the next reopen candidate should be searched first inside the surviving normalized pinning chain.
\end{proposition}

\begin{proof}
By Proposition~\ref{prop:conservative-monte-carlo-qam-probe-rule}, the recurrence class \(\mathcal C\) may only prioritize a route already present in the normalized theorem/tool layer. Hence the total-channel baseline must remain frozen at \(\mathrm{TR}_4\), exactly as in the bridge/localization decoder.

Now assume \(\mathcal C\) is purely local. By definition, its repeated slot/rank pattern already becomes coherent under a local imbalance-side rereading, and no earlier matrix/phase reinterpretation yields a smaller admissible menu. Therefore the decoder is not allowed to jump first to \(M^{\mathrm{surf}}_{\mathrm{br}}\) or \(M^{\mathrm{loc}}_{\mathrm{br}}\). The smallest normalized reopening surface that can already account for the stable recurrence class is instead the surviving pinning menu \(M^{\mathrm{pin}}_{\mathrm{imb}}\).

Inside that menu the current live OP~1 burden has already been compressed in Proposition~\ref{prop:first-tag-first-op1-route}. The first admissible entry is therefore \(\mathrm{IP}_{10}\), the optimum-pinning comparison-rigidity bridge. Once that bridge is opened, one proceeds to \(\mathrm{IP}_9\), the retained-window-rigidity / no-second-carrier / strict-upper-gap bridge, and then to \(\mathrm{IP}_8\), the admissible-window uniqueness bridge. Thus the displayed tag route is exactly the smallest normalized decoding of a stable purely local recurrence class.
\end{proof}

\begin{remark}[What the first Monte-Carlo-to-IP decoding now buys]
\label{rem:what-the-first-monte-carlo-to-ip-decoding-buys}
The decoder is now symmetric across the two live diagnostic sides. The present line may state in compact normalized form: ``freeze at \(\mathrm{TR}_4\); if the recurrence is translation-sensitive, decode first by \(\mathrm{BS}_1\to\mathrm{BS}_3\to\mathrm{BL}_2\); if the recurrence is already imbalance-stable, decode instead by \(\mathrm{IP}_{10}\to\mathrm{IP}_9\to\mathrm{IP}_8\).'' In that precise sense the Monte-Carlo layer no longer only points vaguely toward ``some residual side''; it now separates, conservatively and in dereferenceable tags, the first bridge reopening from the first compressed OP~1 reopening.
\end{remark}


\subsection*{19AZ$^{(4c)}$. First explicit decision table for residual-signature routing}
\label{sec:first-explicit-decision-table-residual-signature-routing}
The decoder layer should now be made directly usable. The present line no longer needs only two separate propositions saying when a stable recurrence class goes to the bridge/localization side or stays on the pinning side. It now needs one explicit decision table that turns the first observed residual-signature class into a smallest admissible action. In particular, the table should also record the conservative ``hold'' outcome when no stable recurrence class appears at all.

\begin{quote}\small
\noindent\textbf{Reader cue.} Read the next table as an action selector. Its purpose is to turn an observed residual-signature class into the smallest admissible reopen instruction, not to claim that every downstream route has already been exhausted.
\end{quote}

\begin{definition}[First explicit decision table for residual-signature routing]
\label{def:first-explicit-decision-table-residual-signature-routing}
Assume the retained total-channel baseline remains frozen through \(\mathrm{TR}_4\), and let \(\mathcal C\) be the first sampled recurrence class produced by the Kasiski-style scan from Definition~\ref{def:kasiski-style-recurrence-scan}. The \emph{first explicit decision table for residual-signature routing} is the map
\[
\mathsf{RouteTab}(\mathcal C)\in \bigl\{\mathsf{Hold},\,\mathsf{IP\text{-}first},\,\mathsf{BS/BL\text{-}first},\,\mathsf{BL\text{-}continue},\,\mathsf{BL\text{-}kernel}\bigr\}
\]
with the following normal form.

\begin{center}
\renewcommand{\arraystretch}{1.2}
\begin{tabular}{>{\raggedright\arraybackslash}p{0.26\linewidth} >{\raggedright\arraybackslash}p{0.22\linewidth} >{\raggedright\arraybackslash}p{0.38\linewidth}}
\toprule
Observed residual-signature class & Decoder output & First normalized route \\
\midrule
No stable recurrence class appears & \(\mathsf{Hold}\) & no reopen; keep the theorem/tool layer frozen \\
Purely local / imbalance-stable recurrence & \(\mathsf{IP\text{-}first}\) & \(\mathrm{TR}_4 \Rightarrow \mathrm{IP}_{10} \Rightarrow \mathrm{IP}_9 \Rightarrow \mathrm{IP}_8\) \\
Translation-sensitive recurrence & \(\mathsf{BS/BL\text{-}first}\) & \(\mathrm{TR}_4 \Rightarrow \mathrm{BS}_1 \Rightarrow \mathrm{BS}_3 \Rightarrow \mathrm{BL}_2\) \\
Translation-sensitive recurrence still live after \(\mathrm{BL}_2\) & \(\mathsf{BL\text{-}continue}\) & append \(\mathrm{BL}_4\) \\
First localized obstruction lies in the readout/kernel stratum & \(\mathsf{BL\text{-}kernel}\) & append \(\mathrm{BL}_5\) \\
\bottomrule
\end{tabular}
\end{center}

The table is read top-down and always returns the smallest normalized action whose recorded carriers already account for the observed recurrence behavior.
\end{definition}

\begin{proposition}[The first explicit decision table is conservative and minimal on the current live classes]
\label{prop:first-explicit-decision-table-conservative-and-minimal}
Under the current theorem/tool layer, every sampled residual-signature class that is presently admitted by Propositions~\ref{prop:translation-sensitive-signature-classes-prioritize-bs-bl-route} and \ref{prop:purely-local-signature-classes-prioritize-ip-route} receives from Definition~\ref{def:first-explicit-decision-table-residual-signature-routing} a unique first admissible action. More precisely:
\begin{enumerate}[label=(\roman*)]
    \item if no stable recurrence class appears, the decision table returns \(\mathsf{Hold}\);
    \item if the recurrence class is purely local, it returns \(\mathsf{IP\text{-}first}\);
    \item if the recurrence class is translation-sensitive, it returns \(\mathsf{BS/BL\text{-}first}\);
    \item deeper localization steps \(\mathsf{BL\text{-}continue}\) and \(\mathsf{BL\text{-}kernel}\) may be appended only after the route has already entered the BS/BL side through \(\mathrm{BL}_2\).
\end{enumerate}
Hence the table is conservative, because it never manufactures a new carrier or bypasses the frozen total-channel anchor, and minimal, because it always picks the smallest currently recorded reopening route compatible with the observed recurrence class.
\end{proposition}

\begin{proof}
If no stable recurrence class appears at all, then Remark~\ref{rem:immediate-diagnostic-reading-signatures} already says that the theorem layer must remain untouched; this is exactly the output \(\mathsf{Hold}\). If a stable recurrence class is purely local, then Proposition~\ref{prop:purely-local-signature-classes-prioritize-ip-route} forces the first admissible reopening route to stay inside the pinning chain beneath \(\mathrm{TR}_4\), so the table must return \(\mathsf{IP\text{-}first}\) and the route \(\mathrm{TR}_4 \Rightarrow \mathrm{IP}_{10} \Rightarrow \mathrm{IP}_9 \Rightarrow \mathrm{IP}_8\). If instead the stable recurrence class is translation-sensitive, then Proposition~\ref{prop:translation-sensitive-signature-classes-prioritize-bs-bl-route} forces the first admissible reopening route to pass through \(\mathrm{BS}_1 \Rightarrow \mathrm{BS}_3 \Rightarrow \mathrm{BL}_2\), so the table must return \(\mathsf{BS/BL\text{-}first}\).

The remaining two outputs are not independent first moves. By Proposition~\ref{prop:translation-sensitive-signature-classes-prioritize-bs-bl-route}, \(\mathrm{BL}_4\) is appended only if the obstruction remains live after \(\mathrm{BL}_2\), and \(\mathrm{BL}_5\) is appended only when that localized obstruction lies in the readout/kernel stratum. Therefore the decision table never jumps directly to \(\mathrm{BL}_4\) or \(\mathrm{BL}_5\) without first exhausting the smaller admissible bridge/localization entry route. This proves both conservativity and minimality on the currently live recurrence classes.
\end{proof}

\begin{remark}[Immediate operational reading of the decision table]
\label{rem:immediate-operational-reading-of-the-decision-table}
The present line can now state the decoder in one sentence that is usable both by a reader and by the future machine layer: classify the first stable residual-signature class, then apply
\[
\mathsf{Hold}\;\vee\;\mathsf{IP\text{-}first}\;\vee\;\mathsf{BS/BL\text{-}first},
\]
with deeper localization allowed only on the BS/BL side by appending \(\mathsf{BL\text{-}continue}\) and, in the readout/kernel case, \(\mathsf{BL\text{-}kernel}\). In that sense the Monte-Carlo diagnostics, the QAM tag layer, and the normalized reopening routes now meet in one explicit working table rather than only in scattered prose.
\end{remark}


\subsection*{19AZ$^{(4d)}$. First deterministic QAM decoder syntax for residual-signature routing}
\label{sec:first-deterministic-qam-decoder-syntax-residual-signature-routing}
The decision table may now be compressed one step further. The present line no longer needs only a prose table saying which class goes to hold, IP-first, or BS/BL-first. It may now record one short deterministic decoder syntax whose words already match the normalized QAM routes. The point is not to introduce a new theory of automata; it is only to give the future machine layer and the human reader the same finite, dereferenceable instruction string.

\begin{definition}[First deterministic QAM decoder syntax]
\label{def:first-deterministic-qam-decoder-syntax}
Assume the retained total-channel anchor has been frozen through \(\mathrm{TR}_4\), and let \(\mathcal C\) be the first sampled residual-signature class from the Kasiski-style recurrence scan. The \emph{first deterministic QAM decoder syntax} is the finite rewrite alphabet
\[
\mathfrak A_{\mathrm{dec}}=\bigl\{\mathsf{classify},\,\mathsf{hold},\,\mathsf{ip},\,\mathsf{bsbl},\,\mathsf{continue},\,\mathsf{kernel}\bigr\}
\]
with the following admissible output words:
\[
\begin{aligned}
\mathsf{classify}&\to\mathsf{hold},\\
\mathsf{classify}&\to\mathsf{ip},\\
\mathsf{classify}&\to\mathsf{bsbl},\\
\mathsf{classify}&\to\mathsf{bsbl}\to\mathsf{continue},\\
\mathsf{classify}&\to\mathsf{bsbl}\to\mathsf{continue}\to\mathsf{kernel}.
\end{aligned}
\]
These words are interpreted by the normalized dereferencing map
\[
\mathsf{Dec}_{\mathrm{QAM}}:
\begin{cases}
\mathsf{hold} &\mapsto \text{keep the theorem/tool layer frozen},\\[0.3em]
\mathsf{ip} &\mapsto \mathrm{TR}_4 \Rightarrow \mathrm{IP}_{10} \Rightarrow \mathrm{IP}_9 \Rightarrow \mathrm{IP}_8,\\[0.3em]
\mathsf{bsbl} &\mapsto \mathrm{TR}_4 \Rightarrow \mathrm{BS}_1 \Rightarrow \mathrm{BS}_3 \Rightarrow \mathrm{BL}_2,\\[0.3em]
\mathsf{continue} &\mapsto \text{append }\mathrm{BL}_4,\\[0.3em]
\mathsf{kernel} &\mapsto \text{append }\mathrm{BL}_5.
\end{cases}
\]
The syntax is called \emph{deterministic} when exactly one of the following first outputs is assigned to \(\mathcal C\): \(\mathsf{hold}\), \(\mathsf{ip}\), or \(\mathsf{bsbl}\); the symbols \(\mathsf{continue}\) and \(\mathsf{kernel}\) may occur only after \(\mathsf{bsbl}\).
\end{definition}

\begin{proposition}[The deterministic decoder syntax is equivalent to the current decision table]
\label{prop:deterministic-decoder-syntax-equivalent-current-decision-table}
Under the current live theorem/tool layer, Definition~\ref{def:first-deterministic-qam-decoder-syntax} yields the same first admissible routing data as Definition~\ref{def:first-explicit-decision-table-residual-signature-routing}. More precisely:
\begin{enumerate}[label=(\roman*)]
    \item if no stable recurrence class appears, the unique admissible output word is \(\mathsf{classify}\to\mathsf{hold}\);
    \item if the recurrence class is purely local / imbalance-stable, the unique admissible output word is \(\mathsf{classify}\to\mathsf{ip}\);
    \item if the recurrence class is translation-sensitive, the unique admissible output word is \(\mathsf{classify}\to\mathsf{bsbl}\);
    \item if, after entering the BS/BL side, the obstruction remains live beyond \(\mathrm{BL}_2\), the decoder extends uniquely to \(\mathsf{classify}\to\mathsf{bsbl}\to\mathsf{continue}\);
    \item if that localized obstruction lies in the readout/kernel stratum, the decoder extends further uniquely to \(\mathsf{classify}\to\mathsf{bsbl}\to\mathsf{continue}\to\mathsf{kernel}\).
\end{enumerate}
Hence the syntax is a faithful finite normal-form compression of the decision table, with no new carriers and no loss of routing information.
\end{proposition}

\begin{proof}
The first three clauses are exactly the three primary outputs of Definition~\ref{def:first-explicit-decision-table-residual-signature-routing}: \(\mathsf{Hold}\), \(\mathsf{IP\text{-}first}\), and \(\mathsf{BS/BL\text{-}first}\). Definition~\ref{def:first-deterministic-qam-decoder-syntax} merely rewrites these three outcomes as the shorter syntax outputs \(\mathsf{hold}\), \(\mathsf{ip}\), and \(\mathsf{bsbl}\), while the dereferencing map \(\mathsf{Dec}_{\mathrm{QAM}}\) sends them back to the same normalized routes.

The remaining two clauses are likewise only compressed spellings of the already recorded BS/BL-side continuations. By Proposition~\ref{prop:first-explicit-decision-table-conservative-and-minimal}, the symbol \(\mathsf{BL\text{-}continue}\) may appear only after the bridge/localization route has already entered through \(\mathrm{BL}_2\), and \(\mathsf{BL\text{-}kernel}\) may appear only after that continuation has localized the first remaining obstruction inside the readout/kernel stratum. In the deterministic syntax these two later actions are rewritten as \(\mathsf{continue}\) and \(\mathsf{kernel}\), with the same ordering restriction. Thus every admissible table output produces exactly one admissible decoder word, and every admissible decoder word dereferences to exactly one table-consistent reopening route.

Therefore the syntax introduces no new mathematical burden. It only compresses the already normalized routing data into a finite QAM-readable word system. This proves equivalence to the current decision table and faithfulness of the compression.
\end{proof}

\begin{remark}[Immediate machine reading of the deterministic decoder syntax]
\label{rem:immediate-machine-reading-deterministic-decoder-syntax}
The present line may now be summarized by the short machine-readable sentences
\[
\mathsf{classify}\to\mathsf{hold},\qquad
\mathsf{classify}\to\mathsf{ip},\qquad
\mathsf{classify}\to\mathsf{bsbl}\to\mathsf{continue}\to\mathsf{kernel},
\]
with the understanding that later symbols may only be appended when the earlier route has actually been entered and justified. In that sense the QAM layer now has not only tags and tables, but also a first deterministic decoder syntax that a future theorem-store, proof assistant, or simulation layer could read without reopening the entire surrounding prose.
\end{remark}



\subsection*{19AZ$^{(4e)}$. First finite transition function for deterministic residual-signature routing}
\label{sec:first-finite-transition-function-deterministic-residual-signature-routing}
The deterministic decoder syntax may now be compressed once more into a tiny state transition rule. The point is again modest: not to introduce a new autonomous automaton theory, but only to expose the present routing discipline as a finite partial transition function that the future machine layer can read directly. In particular, the current line should make explicit that only three first moves are admitted from the classification state, while deeper continuation steps remain locked behind the BS/BL side.

\begin{definition}[First finite transition function for the deterministic decoder]
\label{def:first-finite-transition-function-deterministic-decoder}
Let
\[
Q_{\mathrm{dec}}=\bigl\{\mathsf{classify},\,\mathsf{hold},\,\mathsf{ip},\,\mathsf{bsbl},\,\mathsf{continue},\,\mathsf{kernel}\bigr\}
\]
be the current decoder-state set, and let
\[
\Sigma_{\mathrm{dec}}=\bigl\{\mathsf{none},\,\mathsf{local},\,\mathsf{trans},\,\mathsf{stilllive},\,\mathsf{kernelstratum}\bigr\}
\]
be the present event alphabet, where:
\begin{itemize}[leftmargin=2em]
    \item \(\mathsf{none}\) means that no stable residual-signature recurrence class appears,
    \item \(\mathsf{local}\) means that the sampled class is purely local / imbalance-stable,
    \item \(\mathsf{trans}\) means that the sampled class is translation-sensitive,
    \item \(\mathsf{stilllive}\) means that after entering the BS/BL side the obstruction remains live beyond \(\mathrm{BL}_2\),
    \item \(\mathsf{kernelstratum}\) means that the first localized remaining obstruction lies in the readout/kernel stratum.
\end{itemize}
The \emph{first finite transition function for deterministic residual-signature routing} is the partial map
\[
\delta_{\mathrm{dec}}:Q_{\mathrm{dec}}\times \Sigma_{\mathrm{dec}}\dashrightarrow Q_{\mathrm{dec}}
\]
given by the only admissible transitions
\[
\begin{aligned}
\delta_{\mathrm{dec}}(\mathsf{classify},\mathsf{none})&=\mathsf{hold},\\
\delta_{\mathrm{dec}}(\mathsf{classify},\mathsf{local})&=\mathsf{ip},\\
\delta_{\mathrm{dec}}(\mathsf{classify},\mathsf{trans})&=\mathsf{bsbl},\\
\delta_{\mathrm{dec}}(\mathsf{bsbl},\mathsf{stilllive})&=\mathsf{continue},\\
\delta_{\mathrm{dec}}(\mathsf{continue},\mathsf{kernelstratum})&=\mathsf{kernel}.
\end{aligned}
\]
All other transitions are undefined at the present stage. The state outputs are dereferenced through the same normalized map \(\mathsf{Dec}_{\mathrm{QAM}}\) from Definition~\ref{def:first-deterministic-qam-decoder-syntax}.
\end{definition}

\begin{proposition}[The finite transition function is equivalent to the deterministic decoder syntax and remains minimal on the current live events]
\label{prop:finite-transition-function-equivalent-deterministic-decoder-syntax}
Under the current theorem/tool layer, the partial transition function from Definition~\ref{def:first-finite-transition-function-deterministic-decoder} carries exactly the same routing information as the deterministic decoder syntax from Definition~\ref{def:first-deterministic-qam-decoder-syntax}. More precisely:
\begin{enumerate}[label=(\roman*)]
    \item the three first moves from \(\mathsf{classify}\) are exactly \(\mathsf{hold}\), \(\mathsf{ip}\), and \(\mathsf{bsbl}\), corresponding respectively to the events \(\mathsf{none}\), \(\mathsf{local}\), and \(\mathsf{trans}\);
    \item the only admissible continuation on the bridge/localization side is \(\mathsf{bsbl}\xrightarrow{\mathsf{stilllive}}\mathsf{continue}\);
    \item the only admissible further refinement is \(\mathsf{continue}\xrightarrow{\mathsf{kernelstratum}}\mathsf{kernel}\);
    \item every undefined transition would either skip a smaller admissible reopening route or manufacture a carrier/event type not yet recorded in the present line.
\end{enumerate}
Hence \(\delta_{\mathrm{dec}}\) is a faithful finite-state compression of the current routing discipline and is minimal on the currently live diagnostic events.
\end{proposition}

\begin{proof}
The first three transitions are exactly the three initial output words from Proposition~\ref{prop:deterministic-decoder-syntax-equivalent-current-decision-table}: \(\mathsf{classify}\to\mathsf{hold}\), \(\mathsf{classify}\to\mathsf{ip}\), and \(\mathsf{classify}\to\mathsf{bsbl}\). The transition function merely rewrites these words as event-labeled state moves. The event names \(\mathsf{none}\), \(\mathsf{local}\), and \(\mathsf{trans}\) are exactly the presently live coarse classes isolated by the decision table and the two decoder propositions.

Once the route has entered \(\mathsf{bsbl}\), the deterministic syntax allows only the optional later symbols \(\mathsf{continue}\) and \(\mathsf{kernel}\). By Definition~\ref{def:first-deterministic-qam-decoder-syntax}, \(\mathsf{continue}\) may appear only after \(\mathsf{bsbl}\), and by Proposition~\ref{prop:first-explicit-decision-table-conservative-and-minimal} it is justified precisely when the obstruction remains live beyond \(\mathrm{BL}_2\). This is the transition \(\delta_{\mathrm{dec}}(\mathsf{bsbl},\mathsf{stilllive})=\mathsf{continue}\). Likewise, \(\mathsf{kernel}\) may occur only after that continuation has localized the first remaining obstruction inside the readout/kernel stratum, giving \(\delta_{\mathrm{dec}}(\mathsf{continue},\mathsf{kernelstratum})=\mathsf{kernel}\).

No further transitions are admissible on the current line. Any direct jump from \(\mathsf{classify}\) to \(\mathsf{continue}\) or \(\mathsf{kernel}\) would bypass the smaller bridge/localization entry route, while any jump from \(\mathsf{ip}\) to a later BS/BL state would conflate two presently separated reopening disciplines. And any additional event symbol would exceed the currently normalized recurrence classes. Therefore the partial transition function is equivalent to the deterministic decoder syntax, conservative with respect to the recorded carriers, and minimal on the currently live events.
\end{proof}

\begin{remark}[Immediate automaton reading of the current decoder]
\label{rem:immediate-automaton-reading-current-decoder}
The present line can now be read as a tiny partial automaton:
\[
\begin{aligned}
\mathsf{classify} &\xrightarrow{\mathsf{none}} \mathsf{hold},\\
\mathsf{classify} &\xrightarrow{\mathsf{local}} \mathsf{ip},\\
\mathsf{classify} &\xrightarrow{\mathsf{trans}} \mathsf{bsbl}
\xrightarrow{\mathsf{stilllive}} \mathsf{continue}
\xrightarrow{\mathsf{kernelstratum}} \mathsf{kernel}.
\end{aligned}
\]
Through the dereferencing map \(\mathsf{Dec}_{\mathrm{QAM}}\), this automaton is still nothing more than the already normalized theorem/tool layer. But it now has a form that a later theorem store, proof assistant, or Monte-Carlo controller could read as a deterministic transition discipline instead of as prose plus scattered tags.
\end{remark}


\subsection*{19AZ$^{(4f)}$. First observer/event normal form for deterministic residual-signature words}
\label{sec:first-observer-event-normal-form-deterministic-residual-signature-words}
The tiny decoder automaton from the previous subsection still presupposes that the incoming event has already been named as one of the symbols \(\mathsf{none}\), \(\mathsf{local}\), \(\mathsf{trans}\), \(\mathsf{stilllive}\), or \(\mathsf{kernelstratum}\). The next conservative step is therefore to normalize not the output route once more, but the observer-side input itself. The present line does not need a large language for that purpose. It only needs a short finite event-word normal form recording how much of the bridge/localization continuation has actually become visible at readout.

\begin{definition}[First observer/event normal form for deterministic residual-signature words]
\label{def:first-observer-event-normal-form-deterministic-residual-signature-words}
Let
\[
\mathcal W^{\mathrm{nf}}_{\mathrm{dec}}
=
\bigl\{
\mathsf{none},\,
\mathsf{local},\,
\mathsf{trans},\,
\mathsf{trans}\,\mathsf{stilllive},\,
\mathsf{trans}\,\mathsf{stilllive}\,\mathsf{kernelstratum}
\bigr\}
\]
be the \emph{first observer/event normal-form language} for deterministic residual-signature routing. The associated partial extended transition map is defined recursively by
\[
\delta^{\ast}_{\mathrm{dec}}(q,\varepsilon)=q,
\qquad
\delta^{\ast}_{\mathrm{dec}}(q,aw)=
\delta^{\ast}_{\mathrm{dec}}\bigl(\delta_{\mathrm{dec}}(q,a),w\bigr)
\]
whenever the right-hand side is defined, where \(a\in\Sigma_{\mathrm{dec}}\) and \(w\) is a finite word in the same alphabet.

The \emph{observer/event normal form map}
\[
\mathsf{NF}_{\mathrm{obs}}: \text{currently normalized residual-signature cases}
\longrightarrow \mathcal W^{\mathrm{nf}}_{\mathrm{dec}}
\]
assigns the unique normal-form word as follows:
\begin{enumerate}[label=(\roman*)]
    \item no stable recurrence class is observed \(\mapsto \mathsf{none}\),
    \item a purely local / imbalance-stable class is observed \(\mapsto \mathsf{local}\),
    \item a translation-sensitive class is observed but no remaining obstruction survives beyond \(\mathrm{BL}_2\) \(\mapsto \mathsf{trans}\),
    \item a translation-sensitive class is observed and the obstruction remains live beyond \(\mathrm{BL}_2\) \(\mapsto \mathsf{trans}\,\mathsf{stilllive}\),
    \item the previous case holds and the first localized remaining obstruction lies in the readout/kernel stratum \(\mapsto \mathsf{trans}\,\mathsf{stilllive}\,\mathsf{kernelstratum}\).
\end{enumerate}
\end{definition}

\begin{proposition}[The observer/event normal forms are exactly the currently admissible decoder words and reproduce the normalized reopening routes]
\label{prop:observer-event-normal-forms-exactly-currently-admissible-decoder-words}
Under the current theorem/tool layer, the language \(\mathcal W^{\mathrm{nf}}_{\mathrm{dec}}\) from Definition~\ref{def:first-observer-event-normal-form-deterministic-residual-signature-words} is exactly the set of currently admissible observer-side event words for the deterministic decoder. More precisely:
\begin{enumerate}[label=(\roman*)]
    \item every word in \(\mathcal W^{\mathrm{nf}}_{\mathrm{dec}}\) is admitted by the partial automaton \(\delta_{\mathrm{dec}}\) when started at \(\mathsf{classify}\),
    \item the terminal states are respectively
    \[
    \delta^{\ast}_{\mathrm{dec}}(\mathsf{classify},\mathsf{none})=\mathsf{hold},\qquad
    \delta^{\ast}_{\mathrm{dec}}(\mathsf{classify},\mathsf{local})=\mathsf{ip},
    \]
    \[
    \delta^{\ast}_{\mathrm{dec}}(\mathsf{classify},\mathsf{trans})=\mathsf{bsbl},\qquad
    \delta^{\ast}_{\mathrm{dec}}(\mathsf{classify},\mathsf{trans}\,\mathsf{stilllive})=\mathsf{continue},
    \]
    \[
    \delta^{\ast}_{\mathrm{dec}}(\mathsf{classify},\mathsf{trans}\,\mathsf{stilllive}\,\mathsf{kernelstratum})=\mathsf{kernel},
    \]
    and these states dereference exactly to the normalized routes already fixed in the QAM tool layer;
    \item no other finite event word is currently admissible without either violating the transition discipline of \(\delta_{\mathrm{dec}}\) or introducing an observer-side distinction not yet normalized in the present line.
\end{enumerate}
Hence the current Monte-Carlo/QAM interface may now be read not only as a state machine, but as a finite normal-form language of observer-side diagnostic words.
\end{proposition}

\begin{proof}
The five displayed words are admitted directly by the transition function from Definition~\ref{def:first-finite-transition-function-deterministic-decoder}. The one-letter words \(\mathsf{none}\), \(\mathsf{local}\), and \(\mathsf{trans}\) are exactly the three first-step events recorded there. Once the route has entered \(\mathsf{bsbl}\), the only currently permitted continuation event is \(\mathsf{stilllive}\), and once the route has entered \(\mathsf{continue}\), the only currently permitted further refinement is \(\mathsf{kernelstratum}\). Therefore the longer words \(\mathsf{trans}\,\mathsf{stilllive}\) and \(\mathsf{trans}\,\mathsf{stilllive}\,\mathsf{kernelstratum}\) are admitted as well.

The recursive definition of \(\delta^{\ast}_{\mathrm{dec}}\) then gives the terminal states by straightforward evaluation. For \(\mathsf{none}\) and \(\mathsf{local}\), only one transition is applied, yielding \(\mathsf{hold}\) and \(\mathsf{ip}\). For \(\mathsf{trans}\), the first step yields \(\mathsf{bsbl}\). Appending \(\mathsf{stilllive}\) moves to \(\mathsf{continue}\), and appending \(\mathsf{kernelstratum}\) moves from there to \(\mathsf{kernel}\). By Definition~\ref{def:first-deterministic-qam-decoder-syntax}, each terminal state dereferences to the already normalized reopening route: \(\mathsf{hold}\) means reopen nothing, \(\mathsf{ip}\) means the chain \(\mathrm{TR}_4\Rightarrow \mathrm{IP}_{10}\Rightarrow \mathrm{IP}_9\Rightarrow \mathrm{IP}_8\), \(\mathsf{bsbl}\) means the chain \(\mathrm{TR}_4\Rightarrow \mathrm{BS}_1\Rightarrow \mathrm{BS}_3\Rightarrow \mathrm{BL}_2\), \(\mathsf{continue}\) appends \(\mathrm{BL}_4\), and \(\mathsf{kernel}\) appends \(\mathrm{BL}_5\).

It remains to exclude any additional finite event word. A word beginning with anything other than \(\mathsf{none}\), \(\mathsf{local}\), or \(\mathsf{trans}\) is impossible because \(\delta_{\mathrm{dec}}\) is undefined at \(\mathsf{classify}\) on all other events. A word beginning with \(\mathsf{none}\) or \(\mathsf{local}\) cannot be extended further without leaving the currently normalized routing discipline. And a word beginning with \(\mathsf{trans}\) may only continue with \(\mathsf{stilllive}\), then optionally with \(\mathsf{kernelstratum}\), because no further transitions are currently recorded. Thus every admissible observer-side case has exactly one normal-form word and no other finite word is currently allowed.
\end{proof}

\begin{remark}[Immediate machine reading of the observer/event normal form]
\label{rem:immediate-machine-reading-observer-event-normal-form}
The current live interface may now be summarized by the finite observer-language
\[
\mathsf{none},\qquad
\mathsf{local},\qquad
\mathsf{trans},\qquad
\mathsf{trans}\,\mathsf{stilllive},\qquad
\mathsf{trans}\,\mathsf{stilllive}\,\mathsf{kernelstratum}.
\]
This is the first point at which the tool layer, the QAM tags, the Monte-Carlo recurrence classes, the deterministic decoder syntax, and the partial transition automaton all meet in one compact object: a finite word language on the observer side whose outputs are already normalized reopening routes on the theorem side.
\end{remark}



\subsection*{19AZ$^{(4g)}$. First observer-word evaluation map into deterministic decoder outputs}
\label{sec:first-observer-word-evaluation-map-deterministic-decoder-outputs}
The observer/event normal forms now provide a finite input language for the deterministic decoder. The next conservative compression step is to package the whole passage from normalized observer word to normalized theorem-side reopening output into one explicit evaluation map. This adds no new event class and no new route. It merely composes the already stabilized normal-form language, transition automaton, and QAM dereferencing rule into one finite readout operator.

\begin{definition}[First observer-word evaluation map into deterministic decoder outputs]
\label{def:first-observer-word-evaluation-map-deterministic-decoder-outputs}
Let
\[
\mathcal W^{\mathrm{nf}}_{\mathrm{dec}}
=
\bigl\{
\mathsf{none},\,
\mathsf{local},\,
\mathsf{trans},\,
\mathsf{trans}\,\mathsf{stilllive},\,
\mathsf{trans}\,\mathsf{stilllive}\,\mathsf{kernelstratum}
\bigr\}
\]
be the observer/event normal-form language from Definition~\ref{def:first-observer-event-normal-form-deterministic-residual-signature-words}. Define the \emph{observer-word evaluation map}
\[
\mathsf{Eval}_{\mathrm{obs}}:
\mathcal W^{\mathrm{nf}}_{\mathrm{dec}}
\longrightarrow
Q_{\mathrm{dec}}\times \mathcal R_{\mathrm{tool}}^{\mathrm{out}}
\]
by
\[
\mathsf{Eval}_{\mathrm{obs}}(w)
=
\Bigl(
\delta^{\ast}_{\mathrm{dec}}(\mathsf{classify},w),
\mathsf{Dec}_{\mathrm{QAM}}\bigl(\delta^{\ast}_{\mathrm{dec}}(\mathsf{classify},w)\bigr)
\Bigr),
\]
where \(Q_{\mathrm{dec}}=\{\mathsf{hold},\mathsf{ip},\mathsf{bsbl},\mathsf{continue},\mathsf{kernel}\}\) is the current decoder-state set and \(\mathcal R_{\mathrm{tool}}^{\mathrm{out}}\) denotes the finite set of currently normalized reopening outputs:
\[
\mathcal R_{\mathrm{tool}}^{\mathrm{out}}
=
\Bigl\{
\begin{aligned}
&\varnothing,\\
&[\mathrm{TR}_4,\mathrm{IP}_{10},\mathrm{IP}_9,\mathrm{IP}_8],\\
&[\mathrm{TR}_4,\mathrm{BS}_1,\mathrm{BS}_3,\mathrm{BL}_2],\\
&[\mathrm{TR}_4,\mathrm{BS}_1,\mathrm{BS}_3,\mathrm{BL}_2,\mathrm{BL}_4],\\
&[\mathrm{TR}_4,\mathrm{BS}_1,\mathrm{BS}_3,\mathrm{BL}_2,\mathrm{BL}_4,\mathrm{BL}_5]
\end{aligned}
\Bigr\}.
\]
Equivalently,
\begin{align*}
\mathsf{Eval}_{\mathrm{obs}}(\mathsf{none})
&=(\mathsf{hold},\varnothing),\\
\mathsf{Eval}_{\mathrm{obs}}(\mathsf{local})
&=(\mathsf{ip},[\mathrm{TR}_4,\mathrm{IP}_{10},\mathrm{IP}_9,\mathrm{IP}_8]),\\
\mathsf{Eval}_{\mathrm{obs}}(\mathsf{trans})
&=(\mathsf{bsbl},[\mathrm{TR}_4,\mathrm{BS}_1,\mathrm{BS}_3,\mathrm{BL}_2]),\\
\mathsf{Eval}_{\mathrm{obs}}(\mathsf{trans}\,\mathsf{stilllive})
&=(\mathsf{continue},[\mathrm{TR}_4,\mathrm{BS}_1,\mathrm{BS}_3,\mathrm{BL}_2,\mathrm{BL}_4]),\\
\mathsf{Eval}_{\mathrm{obs}}(\mathsf{trans}\,\mathsf{stilllive}\,\mathsf{kernelstratum})
&=(\mathsf{kernel},[\mathrm{TR}_4,\mathrm{BS}_1,\mathrm{BS}_3,\mathrm{BL}_2,\mathrm{BL}_4,\mathrm{BL}_5]).
\end{align*}
\end{definition}

\begin{proposition}[The observer-word evaluation map is total on the current normal-form language and reproduces the unique normalized reopening outputs]
\label{prop:observer-word-evaluation-map-total-current-normal-form-language}
Under the current theorem/tool layer, the map \(\mathsf{Eval}_{\mathrm{obs}}\) from Definition~\ref{def:first-observer-word-evaluation-map-deterministic-decoder-outputs} is total on \(\mathcal W^{\mathrm{nf}}_{\mathrm{dec}}\). More precisely:
\begin{enumerate}[label=(\roman*)]
    \item every word in \(\mathcal W^{\mathrm{nf}}_{\mathrm{dec}}\) has exactly one decoder terminal state and exactly one currently normalized reopening output under \(\mathsf{Eval}_{\mathrm{obs}}\),
    \item the first component of \(\mathsf{Eval}_{\mathrm{obs}}(w)\) is the unique terminal state reached by the partial automaton \(\delta^{\ast}_{\mathrm{dec}}\),
    \item the second component is exactly the smallest presently admissible theorem-side reopening route compatible with that terminal state,
    \item no currently normalized observer word evaluates to two different reopening outputs.
\end{enumerate}
Hence the present live interface may now be read as one deterministic finite composite map
\[
\mathsf{NF}_{\mathrm{obs}}
\longrightarrow
\mathsf{Eval}_{\mathrm{obs}}
\longrightarrow
\text{normalized reopening output},
\]
rather than as a chain of separately cited decoder pieces.
\end{proposition}

\begin{proof}
By Proposition~\ref{prop:observer-event-normal-forms-exactly-currently-admissible-decoder-words}, every word in \(\mathcal W^{\mathrm{nf}}_{\mathrm{dec}}\) is admitted by the extended transition map \(\delta^{\ast}_{\mathrm{dec}}\) when started at \(\mathsf{classify}\), and its terminal state is uniquely determined. Thus the first component of \(\mathsf{Eval}_{\mathrm{obs}}\) is total and single-valued on the whole normal-form language.

The second component is obtained by applying the deterministic QAM dereferencing map \(\mathsf{Dec}_{\mathrm{QAM}}\) to that unique terminal state. By Definition~\ref{def:first-deterministic-qam-decoder-syntax}, each currently live terminal state carries exactly one normalized reopening output: \(\mathsf{hold}\) gives \(\varnothing\), \(\mathsf{ip}\) gives the fixed IP route \([\mathrm{TR}_4,\mathrm{IP}_{10},\mathrm{IP}_9,\mathrm{IP}_8]\), \(\mathsf{bsbl}\) gives the fixed bridge/localization entry route \([\mathrm{TR}_4,\mathrm{BS}_1,\mathrm{BS}_3,\mathrm{BL}_2]\), \(\mathsf{continue}\) appends \(\mathrm{BL}_4\), and \(\mathsf{kernel}\) appends \(\mathrm{BL}_5\). Therefore the second component is also total and single-valued.

Minimality follows from the construction of the current decoder. The state \(\mathsf{hold}\) reopens nothing, \(\mathsf{ip}\) is already the smallest admissible local route, and the bridge/localization branch enters first through \(\mathsf{bsbl}\) before optional refinement to \(\mathsf{continue}\) and \(\mathsf{kernel}\). Any alternative output would either reopen more than the current state requires or bypass an earlier normalized stage. Hence no current observer word can evaluate to two different normalized reopening outputs.
\end{proof}

\begin{remark}[Immediate operational reading of \texorpdfstring{$\mathsf{Eval}_{\mathrm{obs}}$}{Eval_obs}]
\label{rem:immediate-operational-reading-eval-obs}
The current finite live interface may now be cited in one line:
\[
\begin{aligned}
\text{observer-side recurrence class}
&\xrightarrow{\ \mathsf{NF}_{\mathrm{obs}}\ }
\text{normal-form word}\\
&\xrightarrow{\ \mathsf{Eval}_{\mathrm{obs}}\ }
(\text{decoder state},\text{normalized reopening output}).
\end{aligned}
\]
At this point, the Monte-Carlo diagnostics, the QAM tag layer, the finite decoder automaton, and the theorem-side reopening discipline have become one explicit finite evaluation pipeline. This is the first place in the present line where a later controller could read the whole diagnostic-to-route passage through a single total map on a finite language.
\end{remark}



\subsection*{19AZ$^{(4h)}$. First compact output-normalform tokens for deterministic decoder results}
\label{sec:first-compact-output-normalform-tokens-deterministic-decoder-results}
The observer-word evaluation map now packages every currently admissible observer-side normal form into one of five deterministic theorem-side outputs. The next conservative compression step is to replace those explicit route lists by a finite family of compact output tokens. This adds no new state, no new route, and no new observer word. It merely supplies a shorter QAM-side output carrier for the five already stabilized decoder results.

\begin{definition}[First compact output-normalform tokens for deterministic decoder results]
\label{def:first-compact-output-normalform-tokens-deterministic-decoder-results}
Define the finite set of \emph{output-normalform tokens}
\[
\mathcal T^{\mathrm{nf}}_{\mathrm{out}}
=
\{\mathsf{OH},\mathsf{OI},\mathsf{OB},\mathsf{OC},\mathsf{OK}\},
\]
where the tokens are read as
\begin{align*}
\mathsf{OH}&\leftrightarrow (\mathsf{hold},\varnothing),\\
\mathsf{OI}&\leftrightarrow (\mathsf{ip},[\mathrm{TR}_4,\mathrm{IP}_{10},\mathrm{IP}_9,\mathrm{IP}_8]),\\
\mathsf{OB}&\leftrightarrow (\mathsf{bsbl},[\mathrm{TR}_4,\mathrm{BS}_1,\mathrm{BS}_3,\mathrm{BL}_2]),\\
\mathsf{OC}&\leftrightarrow (\mathsf{continue},[\mathrm{TR}_4,\mathrm{BS}_1,\mathrm{BS}_3,\mathrm{BL}_2,\mathrm{BL}_4]),\\
\mathsf{OK}&\leftrightarrow (\mathsf{kernel},[\mathrm{TR}_4,\mathrm{BS}_1,\mathrm{BS}_3,\mathrm{BL}_2,\mathrm{BL}_4,\mathrm{BL}_5]).
\end{align*}
Equivalently, define the \emph{output-tokenization map}
\[
\mathsf{Tok}_{\mathrm{out}}:
Q_{\mathrm{dec}}\times \mathcal R_{\mathrm{tool}}^{\mathrm{out}}
\longrightarrow
\mathcal T^{\mathrm{nf}}_{\mathrm{out}}
\]
by sending the five currently normalized evaluation outputs from Definition~\ref{def:first-observer-word-evaluation-map-deterministic-decoder-outputs} to the corresponding tokens above. Composing with \(\mathsf{Eval}_{\mathrm{obs}}\) gives the compact output map
\[
\mathsf{Out}^{\mathrm{tok}}_{\mathrm{obs}}
:=
\mathsf{Tok}_{\mathrm{out}}\circ \mathsf{Eval}_{\mathrm{obs}}:
\mathcal W^{\mathrm{nf}}_{\mathrm{dec}}
\longrightarrow
\mathcal T^{\mathrm{nf}}_{\mathrm{out}}.
\]
Thus
\begin{align*}
\mathsf{Out}^{\mathrm{tok}}_{\mathrm{obs}}(\mathsf{none})&=\mathsf{OH},\\
\mathsf{Out}^{\mathrm{tok}}_{\mathrm{obs}}(\mathsf{local})&=\mathsf{OI},\\
\mathsf{Out}^{\mathrm{tok}}_{\mathrm{obs}}(\mathsf{trans})&=\mathsf{OB},\\
\mathsf{Out}^{\mathrm{tok}}_{\mathrm{obs}}(\mathsf{trans}\,\mathsf{stilllive})&=\mathsf{OC},\\
\mathsf{Out}^{\mathrm{tok}}_{\mathrm{obs}}(\mathsf{trans}\,\mathsf{stilllive}\,\mathsf{kernelstratum})&=\mathsf{OK}.
\end{align*}
\end{definition}

\begin{proposition}[The output-normalform tokens are faithful and complete for the current deterministic decoder results]
\label{prop:output-normalform-tokens-faithful-complete-current-deterministic-decoder-results}
Under the current live theorem/tool layer, the token map \(\mathsf{Tok}_{\mathrm{out}}\) is faithful and complete on the presently normalized decoder outputs. More precisely:
\begin{enumerate}[label=(\roman*)]
    \item every value of \(\mathsf{Eval}_{\mathrm{obs}}\) is sent to exactly one token in \(\mathcal T^{\mathrm{nf}}_{\mathrm{out}}\),
    \item different current decoder/output pairs receive different tokens,
    \item the token \(\mathsf{Out}^{\mathrm{tok}}_{\mathrm{obs}}(w)\) determines both the terminal decoder state and the unique normalized reopening route attached to that state,
    \item no information from the current finite live interface is lost when passing from explicit output pairs to output-normalform tokens.
\end{enumerate}
Hence the present diagnostic-to-route pipeline may now also be read in the compact form
\[
\text{observer class}
\longrightarrow
\text{normal-form word}
\longrightarrow
\text{output-normalform token}.
\]
\end{proposition}

\begin{proof}
By Proposition~\ref{prop:observer-word-evaluation-map-total-current-normal-form-language}, the map \(\mathsf{Eval}_{\mathrm{obs}}\) takes values in exactly five currently normalized decoder/output pairs. Definition~\ref{def:first-compact-output-normalform-tokens-deterministic-decoder-results} assigns one token to each of those five pairs and to no others, so every current evaluation value receives exactly one token.

Faithfulness is immediate from the construction: \(\mathsf{OH},\mathsf{OI},\mathsf{OB},\mathsf{OC},\mathsf{OK}\) are pairwise distinct symbols, each tied to one different decoder/output pair. Thus two different current evaluation outputs cannot collapse to the same token. Conversely, each token has exactly one prescribed decoder state and exactly one prescribed normalized reopening route, so the token determines the full current output pair.

Therefore the passage from explicit pairs to compact tokens is lossless on the present finite interface. It is merely a normalization of notation, not a quotient by any nontrivial equivalence. The compact map \(\mathsf{Out}^{\mathrm{tok}}_{\mathrm{obs}}\) preserves exactly the same current routing content as \(\mathsf{Eval}_{\mathrm{obs}}\), only in shorter QAM-readable form.
\end{proof}

\begin{remark}[Why the compact output tokens matter]
\label{rem:why-compact-output-tokens-matter}
The live interface can now be written in its shortest current form:
\[
\text{observer-side recurrence class}
\xrightarrow{\ \mathsf{NF}_{\mathrm{obs}}\ }
\text{normal-form word}
\xrightarrow{\ \mathsf{Out}^{\mathrm{tok}}_{\mathrm{obs}}\ }
\{\mathsf{OH},\mathsf{OI},\mathsf{OB},\mathsf{OC},\mathsf{OK}\}.
\]
At this point the theorem-side reopening discipline already has a compact output alphabet. A later controller or QAM overlay no longer needs to carry the full explicit route list every time; it may cite the token first and dereference the corresponding route only when the detailed theorem-side reopening path is actually needed.
\end{remark}


\subsection*{19AZ$^{(5a)}$. First signature-injectivity normal form for the live OP~1 kernel}
\label{sec:first-signature-injectivity-normal-form-live-op1-kernel}
The normalized routing layer has now done its job on the OP~1 side: it has compressed the live reopening path to
\[
\mathrm{TR}_4 \Longrightarrow \mathrm{IP}_{10} \Longrightarrow \mathrm{IP}_9 \Longrightarrow \mathrm{IP}_8.
\]
The next honest mathematical move is therefore not another routing refinement but a sharper obstruction normal form for the first live carrier in that chain. The remaining burden should now be stated as a local injectivity problem for the retained pinning signature itself: can two genuinely distinct admissible primitive photon-first-zero loops occupy the same retained post-frame corridor class while preserving the same total-channel profile, the same retained optimum location, and the same local comparison/defect readout?

\begin{quote}\small
\noindent\textbf{Reader cue.} Read the next block as an obstruction normal form for the live OP~1 kernel. Its purpose is to rewrite the remaining burden as exclusion of a second admissible primitive loop with the same retained pinning signature, not to claim that this exclusion has already been proved.
\end{quote}

\begin{definition}[Retained pinning signature on the OP~1 corridor]
\label{def:retained-pinning-signature-op1-corridor}
Assume the same-corridor reversal certificate of Definition~\ref{def:op1-same-corridor-reversal}. For every admissible primitive photon-first-zero closed-loop realization $L$ on the retained post-frame corridor, define the \emph{retained pinning signature}
\[
\mathsf{Sig}_{\mathrm{pin}}(L)
\]
to be the complete retained comparison packet carried by the present line on the intended admissible windows, namely: the retained total-channel profile, the retained optimum location, and the standing local comparison/defect data attached to $L$ in the compressed machine-shadow reading. Thus two loops have the same retained pinning signature exactly when they are indistinguishable to the current OP~1 comparison line at the level that feeds Definition~\ref{def:op1-optimum-pinning-comparison-rigidity}.
\end{definition}

\begin{definition}[Duplicate-signature witness for the live OP~1 kernel]
\label{def:duplicate-signature-witness-live-op1-kernel}
Assume Definition~\ref{def:retained-pinning-signature-op1-corridor}. A \emph{duplicate-signature witness} for the live OP~1 kernel is a pair $(L,L')$ of admissible primitive photon-first-zero closed-loop realizations on the retained post-frame corridor such that:
\begin{enumerate}[label=(\roman*)]
    \item $L$ and $L'$ are genuinely distinct up to the standing admissible corridor equivalence,
    \item both preserve the same retained total-channel profile on the intended admissible windows,
    \item \(\mathsf{Sig}_{\mathrm{pin}}(L)=\mathsf{Sig}_{\mathrm{pin}}(L')\).
\end{enumerate}
Equivalently, a duplicate-signature witness is a second primitive loop placement that the current comparison line cannot distinguish from the intended retained loop even though it remains genuinely distinct inside the retained corridor class.
\end{definition}

\begin{proposition}[The live OP~1 burden is equivalent to exclusion of duplicate-signature witnesses]
\label{prop:live-op1-burden-equivalent-exclusion-duplicate-signature-witnesses}
Assume the same-corridor reversal certificate of Definition~\ref{def:op1-same-corridor-reversal}. Then the following are equivalent on the intended admissible windows:
\begin{enumerate}[label=(\roman*)]
    \item the optimum-pinning comparison-rigidity certificate of Definition~\ref{def:op1-optimum-pinning-comparison-rigidity} holds;
    \item no duplicate-signature witness exists in the sense of Definition~\ref{def:duplicate-signature-witness-live-op1-kernel}.
\end{enumerate}
Consequently, together with Proposition~\ref{prop:op1-rigidity-gap-uniqueness-bridge} and the standing no-second-carrier / strict-upper-gap exclusions already built into the current line, exclusion of duplicate-signature witnesses implies
\[
\Delta_{\mathrm{CL}}=0.
\]
Equivalently, the present residual OP~1 burden may now be read exactly as the injectivity of the retained pinning signature on the admissible primitive loop class carried by the retained post-frame corridor.
\end{proposition}

\begin{proof}
Definition~\ref{def:op1-optimum-pinning-comparison-rigidity} says that no genuinely distinct admissible primitive loop on the retained post-frame corridor can preserve the retained total-channel profile, the retained optimum location, and the standing local comparison/defect data while remaining distinct up to the admissible corridor equivalence. But this is exactly the negation of a duplicate-signature witness from Definition~\ref{def:duplicate-signature-witness-live-op1-kernel}, since the retained pinning signature packages precisely those same preserved data. Thus \emph{(i)} and \emph{(ii)} are equivalent.

Once this injectivity statement is available, Proposition~\ref{prop:op1-optimum-pinning-implies-rigidity} yields the retained-window rigidity certificate. Proposition~\ref{prop:op1-rigidity-gap-uniqueness-bridge} then upgrades that rigidity, under the standing no-second-carrier and strict-upper-gap exclusions, to admissible-window uniqueness and hence to \(\Delta_{\mathrm{CL}}=0\).
\end{proof}

\begin{remark}[What the new normal form buys for the next OP~1 step]
\label{rem:what-new-normal-form-buys-next-op1-step}
This is a real sharpening even though it is not yet a closure theorem. Before the present step, the residual burden was phrased as ``prove optimum-pinning comparison rigidity.'' After the present step, the live target is more concrete: show that the retained comparison packet
\[
\mathsf{Sig}_{\mathrm{pin}}
\]
is injective on admissible primitive first-zero loops in the retained post-frame corridor class. In practical terms, the next OP~1 attack should therefore look for a contradiction from a hypothetical duplicate-signature witness, not for a new free continuum-limit parameter. In the normalized QAM reading, this says that the first live carrier \(\mathrm{IP}_{10}\) should now be treated as a signature-injectivity problem for the retained pinning packet.
\end{remark}




\subsection*{19AZ$^{(5b)}$. First contradiction channel split for a duplicate-signature witness}
\label{sec:first-contradiction-channel-split-duplicate-signature-witness}
The normal form of 19AZ$^{(5a)}$ says that the live OP~1 burden is equivalent to the exclusion of duplicate-signature witnesses. The next honest move is therefore to ask what such a witness would have to do inside the present line. Because the retained pinning signature already fixes the retained total-channel profile, the retained optimum location, and the local comparison/defect packet, any residual distinction between two loops carrying the same signature has only two places left to live: either as a second primitive occupancy of the same retained corridor slot, or as a second admissible completion above that same slot. This is exactly the old no-second-carrier / strict-upper-gap dichotomy, but now pulled back to the new signature normal form.

\begin{quote}\small
\noindent\textbf{Reader cue.} Read the next block as the first contradiction routing split for the live OP~1 kernel. Its purpose is to show that a hypothetical duplicate-signature witness can only survive by feeding one of the two exclusion gates already built into the current line: duplicate primitive occupancy or duplicate admissible completion.
\end{quote}

\begin{definition}[Occupancy-side and completion-side duplicate-signature channels]
\label{def:occupancy-completion-channels-duplicate-signature-op1}
Assume the same-corridor reversal certificate of Definition~\ref{def:op1-same-corridor-reversal} and let $(L,L')$ be a duplicate-signature witness in the sense of Definition~\ref{def:duplicate-signature-witness-live-op1-kernel}. We say that the witness is of \emph{occupancy-side type} if the residual distinction between $L$ and $L'$ already appears as a second primitive occupancy of the retained post-frame corridor while the retained total-channel profile is kept fixed. We say that the witness is of \emph{completion-side type} if, after the common retained corridor slot and retained comparison packet are fixed, the residual distinction can only survive as a second admissible completion route above that same retained corridor while preserving the same retained status surface.
\end{definition}

\begin{proposition}[Every duplicate-signature witness feeds a standing OP~1 exclusion gate]
\label{prop:duplicate-signature-witness-feeds-standing-op1-exclusion-gate}
Assume the same-corridor reversal certificate of Definition~\ref{def:op1-same-corridor-reversal}. Let $(L,L')$ be a duplicate-signature witness in the sense of Definition~\ref{def:duplicate-signature-witness-live-op1-kernel}. Then at least one of the following holds on the intended admissible windows:
\begin{enumerate}[label=(\roman*)]
    \item $(L,L')$ is of occupancy-side type in the sense of Definition~\ref{def:occupancy-completion-channels-duplicate-signature-op1};
    \item $(L,L')$ is of completion-side type in the sense of Definition~\ref{def:occupancy-completion-channels-duplicate-signature-op1}.
\end{enumerate}
Consequently, every hypothetical failure of signature-injectivity for the live OP~1 kernel already routes into one of the two standing exclusion channels behind Proposition~\ref{prop:op1-rigidity-gap-uniqueness-bridge}: no-second-carrier on the retained corridor, or strict-upper-gap on admissible completion above that corridor.
\end{proposition}

\begin{proof}
By Definition~\ref{def:duplicate-signature-witness-live-op1-kernel}, the loops $L$ and $L'$ are genuinely distinct up to the standing admissible corridor equivalence, yet they preserve the same retained pinning signature. In particular, they preserve the same retained total-channel profile, the same retained optimum location, and the same local comparison/defect data.

Therefore the current line has already fixed everything that belongs to the retained comparison packet. If the residual distinction between $L$ and $L'$ is visible already at the level of primitive occupancy of the retained post-frame corridor, then the witness is of occupancy-side type by definition. Otherwise the common retained corridor slot and retained comparison packet do not separate the loops, so the remaining distinction can only survive as a second admissible completion above that same retained corridor while preserving the same retained status surface; this is exactly the completion-side case.

Hence every duplicate-signature witness must feed at least one of the two channels. The final sentence is then only a translation into the standing OP~1 language: occupancy-side witnesses are precisely the kind of residual duplication targeted by the no-second-carrier discipline, while completion-side witnesses are precisely the kind of residual duplication targeted by the strict-upper-gap exclusion in Proposition~\ref{prop:op1-rigidity-gap-uniqueness-bridge}.
\end{proof}

\begin{remark}[What this contradiction split buys for the next OP~1 attack]
\label{rem:what-contradiction-split-buys-next-op1-attack}
This still does not close OP~1, but it tells us where a closure attempt must now bite. A hypothetical duplicate-signature witness no longer floats as an abstract rival loop. It is forced into one of two already named contradiction channels. In practical terms, the next OP~1 step should try to show that the retained comparison packet cannot support either occupancy-side duplication on the retained corridor or completion-side duplication above that corridor. In the normalized QAM reading, this means that the live carrier \(\mathrm{IP}_{10}\) has now been resolved into a two-channel contradiction search whose outputs are already aligned with the standing no-second-carrier / strict-upper-gap gates feeding \(\mathrm{IP}_9\).
\end{remark}



\subsection*{19AZ$^{(5c)}$. Occupancy-side duplicate-signature witnesses collapse under the retained no-second-carrier discipline}
\label{sec:occupancy-side-duplicate-signature-witnesses-collapse-no-second-carrier}
After the channel split of 19AZ$^{(5b)}$, the cleanest next question is whether both contradiction channels are still genuinely live. On the present line, the occupancy-side channel sits closest to the already retained theorematic infrastructure: if a duplicate-signature witness already survives as a second primitive occupancy of the same retained corridor while the retained total-channel profile stays fixed, then it is no longer merely a new OP~1 curiosity. It is exactly the sort of duplication that the retained no-second-carrier discipline was built to forbid on the preserved total-channel side.

\begin{quote}\small
\noindent\textbf{Reader cue.} Read the next block as the first actual collapse theorem inside the new OP~1 contradiction split. Its purpose is to show that occupancy-side duplicate-signature witnesses are already excluded by the standing retained no-second-carrier discipline, so that the live burden cannot honestly remain two-channel after this point.
\end{quote}

\begin{proposition}[Occupancy-side duplicate-signature witnesses are excluded on the retained OP~1 corridor]
\label{prop:occupancy-side-duplicate-signature-witnesses-excluded-retained-op1-corridor}
Assume the same-corridor reversal certificate of Definition~\ref{def:op1-same-corridor-reversal}. Assume moreover that the retained no-second-carrier discipline continues to hold on the intended admissible windows at the retained total-channel level, in the compressed sense already available through Theorem~\ref{thm:first-no-second-carrier-theorem-on-the-terminal-interface}. Then no duplicate-signature witness of occupancy-side type exists in the sense of Definitions~\ref{def:duplicate-signature-witness-live-op1-kernel} and~\ref{def:occupancy-completion-channels-duplicate-signature-op1}.
\end{proposition}

\begin{proof}
Suppose for contradiction that \((L,L')\) is a duplicate-signature witness of occupancy-side type. By Definition~\ref{def:duplicate-signature-witness-live-op1-kernel}, the loops \(L\) and \(L'\) are genuinely distinct up to the standing admissible corridor equivalence, yet they preserve the same retained pinning signature. In particular, they preserve the same retained total-channel profile on the intended admissible windows.

By Definition~\ref{def:occupancy-completion-channels-duplicate-signature-op1}, the residual distinction already survives as a second primitive occupancy of the retained post-frame corridor while that retained total-channel profile is kept fixed. But this is precisely the kind of extra primitive retained-carrier occupancy excluded by the standing no-second-carrier discipline at retained total-channel level. Hence the assumed witness contradicts the retained no-second-carrier theorematic layer. Therefore no occupancy-side duplicate-signature witness can exist.
\end{proof}

\begin{corollary}[The live OP~1 duplicate-signature burden reduces to the completion-side channel]
\label{cor:live-op1-duplicate-signature-burden-reduces-completion-side-channel}
Under the hypotheses of Proposition~\ref{prop:occupancy-side-duplicate-signature-witnesses-excluded-retained-op1-corridor}, the following are equivalent on the intended admissible windows:
\begin{enumerate}[label=(\roman*)]
    \item a duplicate-signature witness for the live OP~1 kernel exists;
    \item a duplicate-signature witness of completion-side type exists.
\end{enumerate}
Consequently, after the present step the live OP~1 kernel is no longer honestly read as a two-channel contradiction search. The occupancy-side branch is already discharged by the retained no-second-carrier discipline, so the surviving burden is the exclusion of a second admissible completion above the retained corridor that preserves the same retained pinning signature.
\end{corollary}

\begin{proof}
If a duplicate-signature witness exists, Proposition~\ref{prop:duplicate-signature-witness-feeds-standing-op1-exclusion-gate} shows that it must be either of occupancy-side type or of completion-side type. Proposition~\ref{prop:occupancy-side-duplicate-signature-witnesses-excluded-retained-op1-corridor} eliminates the occupancy-side case. Hence every surviving duplicate-signature witness must be completion-side. The converse is immediate, since every completion-side witness is in particular a duplicate-signature witness.
\end{proof}

\begin{remark}[What this first collapse buys for the next OP~1 step]
\label{rem:what-first-collapse-buys-next-op1-step}
This is the first point in the new post-release OP~1 line where one of the two contradiction channels is not merely named but actually discharged. The live burden is therefore smaller than it looked in 19AZ$^{(5b)}$: the retained comparison packet is already rigid enough against duplicate \emph{occupancy} once the inherited no-second-carrier discipline is invoked. What remains is subtler and now sharply localized. One must show that the same retained pinning signature cannot support a second \emph{admissible completion} above the already fixed retained corridor. In the normalized QAM reading, the live carrier \(\mathrm{IP}_{10}\) has therefore collapsed from a two-channel contradiction search to a completion-side exclusion problem feeding directly into the strict-upper-gap gate behind \(\mathrm{IP}_9\).
\end{remark}




\noindent\textbf{Reader cue.} Read the next block as the second actual collapse theorem in the new OP~1 contradiction split. Its purpose is to show that completion-side duplicate-signature witnesses are already excluded by the standing retained strict-upper-gap discipline, so that no duplicate-signature witness survives the present OP~1 kernel once both standing exclusion gates are read in the new signature normal form.

\subsection*{19AZ$^{(5d)}$. Completion-side duplicate-signature witnesses collapse under the retained strict-upper-gap discipline}
\label{sec:completion-side-duplicate-signature-witnesses-collapse-strict-upper-gap}
After 19AZ$^{(5c)}$, the live duplicate-signature burden has only one honest place left to survive: as a second admissible completion above the already fixed retained corridor while the retained pinning signature stays unchanged. But this is exactly the sort of residual duplication that the retained strict-upper-gap language was already designed to exclude. Once the retained corridor slot and retained comparison packet are fixed, a genuinely second admissible completion preserving the same retained status surface is no longer a new geometric freedom; it is a forbidden duplicate completion above the same retained corridor.

\begin{proposition}[Completion-side duplicate-signature witnesses are excluded on the retained OP~1 corridor]
\label{prop:completion-side-duplicate-signature-witnesses-excluded-retained-op1-corridor}
Assume the same-corridor reversal certificate of Definition~\ref{def:op1-same-corridor-reversal}. Assume moreover that the retained strict-upper-gap discipline continues to hold on the intended admissible windows in the compressed sense already used by Proposition~\ref{prop:op1-rigidity-gap-uniqueness-bridge}. Then no duplicate-signature witness of completion-side type exists in the sense of Definitions~\ref{def:duplicate-signature-witness-live-op1-kernel} and~\ref{def:occupancy-completion-channels-duplicate-signature-op1}.
\end{proposition}

\begin{proof}
Suppose for contradiction that \((L,L')\) is a duplicate-signature witness of completion-side type. By Definition~\ref{def:duplicate-signature-witness-live-op1-kernel}, the loops \(L\) and \(L'\) are genuinely distinct up to the standing admissible corridor equivalence, yet they preserve the same retained pinning signature. In particular, they preserve the same retained total-channel profile, the same retained optimum location, and the same local comparison/defect packet on the intended admissible windows.

By Definition~\ref{def:occupancy-completion-channels-duplicate-signature-op1}, once the common retained corridor slot and retained comparison packet are fixed, the residual distinction between \(L\) and \(L'\) can only survive as a second admissible completion above that same retained corridor while preserving the same retained status surface. But this is precisely the kind of duplicate admissible completion excluded by the standing strict-upper-gap discipline in the retained OP~1 reading. Hence the assumed witness contradicts the retained strict-upper-gap layer. Therefore no completion-side duplicate-signature witness can exist.
\end{proof}

\begin{corollary}[No duplicate-signature witness survives the retained OP~1 kernel]
\label{cor:no-duplicate-signature-witness-survives-retained-op1-kernel}
Assume the hypotheses of Proposition~\ref{prop:occupancy-side-duplicate-signature-witnesses-excluded-retained-op1-corridor} and Proposition~\ref{prop:completion-side-duplicate-signature-witnesses-excluded-retained-op1-corridor}. Then no duplicate-signature witness for the live OP~1 kernel exists on the intended admissible windows.
\end{corollary}

\begin{proof}
If a duplicate-signature witness existed, Corollary~\ref{cor:live-op1-duplicate-signature-burden-reduces-completion-side-channel} would force it to be of completion-side type. Proposition~\ref{prop:completion-side-duplicate-signature-witnesses-excluded-retained-op1-corridor} eliminates that case. Hence no duplicate-signature witness survives.
\end{proof}

\begin{theorem}[Closure of the live OP~1 kernel under the standing retained exclusion disciplines]
\label{thm:closure-live-op1-kernel-standing-retained-exclusion-disciplines}
Assume the same-corridor reversal certificate of Definition~\ref{def:op1-same-corridor-reversal}. Assume moreover that the retained no-second-carrier discipline and the retained strict-upper-gap discipline both hold on the intended admissible windows in the senses used above. Then the optimum-pinning comparison-rigidity certificate of Definition~\ref{def:op1-optimum-pinning-comparison-rigidity} holds on the live OP~1 kernel. Consequently,
\[
\Delta_{\mathrm{CL}}=0.
\]
Equivalently, under the standing retained exclusion disciplines the normalized OP~1 route
\[
\mathrm{TR}_4 \Longrightarrow \mathrm{IP}_{10} \Longrightarrow \mathrm{IP}_9 \Longrightarrow \mathrm{IP}_8
\]
closes on the intended admissible windows.
\end{theorem}

\begin{proof}
By Corollary~\ref{cor:no-duplicate-signature-witness-survives-retained-op1-kernel}, no duplicate-signature witness exists. Proposition~\ref{prop:live-op1-burden-equivalent-exclusion-duplicate-signature-witnesses} therefore yields the optimum-pinning comparison-rigidity certificate. The final sentence of that same proposition then gives \(\Delta_{\mathrm{CL}}=0\) via the already established retained-window-rigidity / admissible-window-uniqueness bridge.
\end{proof}

\begin{remark}[What the completion-side collapse means for the role of the Monte-Carlo probes]
\label{rem:completion-side-collapse-role-monte-carlo-probes}
After the present step, the Monte-Carlo probes with the fusion/decoder reading should not be read as replacements for the OP~1 theorematic closure. Their real leverage is earlier in the search. They compress the search space by classifying residual behavior into a small number of decoder-relevant recurrence classes and by telling us whether a surviving defect should first be treated as \texttt{ip}-local or as \texttt{bsbl}-translation-sensitive. That is why the work now feels faster: the probes do not prove the retained no-second-carrier or strict-upper-gap exclusions, but they stop us from reopening the wrong parts of the document and from carrying broad rival-loop pictures once the residual signature is already locally pinned.

In the present OP~1 line, this means the probes mainly accelerated the route to the right normal form. They helped isolate the live burden as an injectivity problem, then as a two-channel contradiction search, and finally as the completion-side residue that can be handed directly to the strict-upper-gap gate. So the speed-up is real, but it is a \emph{search-space reduction and routing advantage}, not a substitute for the structural closure theorems themselves.
\end{remark}



\noindent\textbf{Reader cue.} Read the next block as the first honest post-OP~1 normalization of the live OP~5 Step~(B) burden. Its purpose is not to declare the routed photon-frame balance defect solved, but to identify the smallest currently licensed structural kernel on which that burden now lives: the production or exclusion of a carrier-upgraded routed-balance witness at the routed photon-frame cutoff.

\subsection*{19AZ$^{(6a)}$. First carrier-upgraded witness normal form for the live OP~5 Step~(B) kernel}
\label{sec:first-carrier-upgraded-witness-normal-form-live-op5b-kernel}
After the closure of the live OP~1 kernel, the honest next burden is OP~5 Step~(B), i.e. the routed photon-frame balance defect \(\Delta_{\Theta R}(T)\). The current line should not reopen that burden through every packet-wave, breathing-envelope, frame-flicker, or helix picture at once. The earlier carrier-first OP~5 architecture already teaches that visible mismatches come in two grades: either they remain screening-only and therefore subordinate to one still-available common closure-bearing return-side carrier package, or they upgrade and certify loss of that common carrier package itself. The cleanest first OP~5(B) normalization is therefore to ask not for a new global formula immediately, but for the smallest carrier-grade witness through which a still-live routed balance defect could honestly persist.

\begin{definition}[Routed-balance witness for the live OP~5 Step~(B) kernel]
\label{def:routed-balance-witness-live-op5b-kernel}
Fix the routed photon-frame cutoff \(T=\gamma_L^{-3/2}\) and the routed balance defect \(\Delta_{\Theta R}(T)\) from Definition~\ref{def:delta-theta-r-residual}. Under the standing carrier-first OP~5 architecture, an admitted visible witness \(w\) is called a \emph{routed-balance witness for the live OP~5 Step~(B) kernel} if \(w\) is cited to support the claim that the routed photon-frame balance layer remains structurally live at the cutoff \(T\). Such a witness is called
\begin{enumerate}[leftmargin=2em]
  \item \emph{screening-only} if it remains screening-only in the sense of Lemma~\ref{lem:carrier-witness-separation};
  \item \emph{carrier-upgraded} if it is carrier-upgraded in the same sense.
\end{enumerate}
Thus a routed-balance witness is not a new witness species; it is the OP~5 Step~(B) reading of the already licensed witness-grade dichotomy at the routed photon-frame cutoff.
\end{definition}

\begin{proposition}[The live OP~5 Step~(B) kernel is a carrier-upgraded witness problem]
\label{prop:live-op5b-kernel-carrier-upgraded-witness-problem}
Assume Proposition~\ref{prop:carrier-first-op5-interface}, Proposition~\ref{prop:packet-phase-epsilon-screening}, Lemma~\ref{lem:first-diagnostic-witness}, Lemma~\ref{lem:screening-witness-compression}, Lemma~\ref{lem:carrier-witness-separation}, and Corollary~\ref{cor:carrier-escalation-threshold}. Then the honest live OP~5 Step~(B) kernel on the present line is the following alternative:
\begin{enumerate}[leftmargin=2em,label=(\roman*)]
  \item either every routed-balance witness at the routed photon-frame cutoff remains screening-only, in which case no carrier/package-level OP~5 Step~(B) obstruction has yet been exhibited;
  \item or there exists a carrier-upgraded routed-balance witness, in which case the live OP~5 Step~(B) burden has genuinely reached the carrier/package layer.
\end{enumerate}
In particular, the first honest structural reopening of \(\Delta_{\Theta R}(T)\) is not an arbitrary visible mismatch, but the production or exclusion of a carrier-upgraded routed-balance witness.
\end{proposition}

\begin{proof}
By Lemma~\ref{lem:carrier-witness-separation}, every admitted visible witness at the carrier-first OP~5 layer lies in exactly one of two grades: screening-only or carrier-upgraded. Definition~\ref{def:routed-balance-witness-live-op5b-kernel} does nothing more than read that existing dichotomy at the routed photon-frame cutoff \(T\) and relative to the residual \(\Delta_{\Theta R}(T)\). If a cited witness remains screening-only, then by Corollary~\ref{cor:carrier-escalation-threshold} its upstream content is still representable through some common closure-bearing return-side carrier package; hence no carrier/package-level OP~5 Step~(B) obstruction has yet been exhibited. If, on the other hand, a cited witness upgrades to carrier grade, then the same corollary says that no such common carrier package remains available, so the routed balance burden has genuinely reached the carrier/package layer. Since no third currently licensed witness grade exists, the displayed alternative is exhaustive. Therefore the first honest structural OP~5 Step~(B) kernel is exactly the production or exclusion of a carrier-upgraded routed-balance witness.
\end{proof}

\begin{corollary}[First normalized reopening route for the live OP~5 Step~(B) kernel]
\label{cor:first-normalized-reopening-route-live-op5b-kernel}
Under the hypotheses of Proposition~\ref{prop:live-op5b-kernel-carrier-upgraded-witness-problem}, the smallest currently licensed normalized reopening route for a genuinely live OP~5 Step~(B) witness is
\[
\mathrm{TR}_4 \Longrightarrow \mathrm{BS}_1 \Longrightarrow \mathrm{BS}_3 \Longrightarrow \mathrm{BL}_2 \Longrightarrow \mathrm{BL}_4 \Longrightarrow \mathrm{BL}_5,
\]
with the last two steps used only when the routed-balance witness remains live after the masked-single-carrier test and localizes at the readout/kernel stratum.
\end{corollary}

\begin{proof}
The total-channel baseline remains frozen by \(\mathrm{TR}_4\). Because the routed photon-frame balance defect is not a pure local optimum-pinning problem but a return-side witness-grade problem, the first honest move is bridge-side rather than imbalance-side: the route enters through the matrix/phase surface-change gates \(\mathrm{BS}_1\) and \(\mathrm{BS}_3\). Once that surface change has been made, the first honest diagnostic question is whether the putative routed-balance witness is still compatible with one masked single carrier, which is precisely the gate \(\mathrm{BL}_2\). If it collapses there, no deeper reopening is licensed. If it remains live, then \(\mathrm{BL}_4\) is exactly the normalized stratification step that localizes the first remaining obstruction. And when that localized obstruction sits in the readout/kernel stratum, \(\mathrm{BL}_5\) is the recorded compression to absence of one common shell-stable carrier package. Therefore the displayed BS/BL route is exactly the smallest currently licensed reopening route for a genuinely live OP~5 Step~(B) witness.
\end{proof}

\begin{remark}[What this first OP~5 Step~(B) normalization buys]
\label{rem:what-first-op5b-normalization-buys}
This is not yet an OP~5 Step~(B) closure theorem. The routed photon-frame balance defect \(\Delta_{\Theta R}(T)\) is still honestly live. But the burden is now sharper. It is no longer carried as an undifferentiated packet/helix/flicker mismatch surface. It is carried as a witness-grade question: does any routed-balance witness actually upgrade to carrier/package level, or do all currently visible routed witnesses remain screening-only beneath one still-available common return-side carrier package? In the present QAM language, OP~1 was compressed onto the \(\mathrm{IP}\)-route, whereas the live OP~5 Step~(B) burden now announces itself as a \(\mathrm{BS}/\mathrm{BL}\)-route burden from the start.
\end{remark}


\noindent\textbf{Reader cue.} Read the next block as the first negative-side sharpening of the live OP~5 Step~(B) kernel. Its purpose is not yet to close the routed photon-frame balance defect, but to identify the first retained situation in which a routed-balance witness cannot honestly escalate to carrier/package grade at all.

\subsection*{19AZ$^{(6b)}$. Carrier-stable routed returns prevent witness-grade escalation on the live OP~5 Step~(B) kernel}
\label{sec:carrier-stable-routed-returns-prevent-witness-grade-escalation-live-op5b-kernel}
The previous normalization reduced the live OP~5 Step~(B) burden to the production or exclusion of a carrier-upgraded routed-balance witness. The next honest question is whether every routed-balance witness must be tested all the way down to the carrier/package layer, or whether one already has a retained positive branch on which such escalation is structurally forbidden. The earlier routed-return package answers this. Once the routed realization lies on the carrier-stable routed return normal form, every admitted visible routed mismatch is already subordinate to one common closure-bearing return-side carrier. In that situation the witness-grade question collapses on the positive side: a routed-balance witness may still exist as screening-level visible data, but it can no longer count as a carrier-grade OP~5 Step~(B) witness.

\begin{proposition}[Carrier-stable routed returns prevent routed-balance witness escalation]
\label{prop:carrier-stable-routed-returns-prevent-routed-balance-witness-escalation}
Assume Proposition~\ref{prop:live-op5b-kernel-carrier-upgraded-witness-problem}, Corollary~\ref{cor:first-normalized-reopening-route-live-op5b-kernel}, Corollary~\ref{cor:carrier-escalation-threshold}, Theorem~\ref{thm:carrier-stable-diagnostic-collapse}, and Theorem~\ref{thm:first-routed-return-normal-form}. Let $(b,r,s;w)$ be a routed branch-shell realization at the cutoff \(T=\gamma_L^{-3/2}\) together with a routed-balance witness \(w\) in the sense of Definition~\ref{def:routed-balance-witness-live-op5b-kernel}. If $(b,r,s;w)$ lies on the \emph{carrier-stable routed return normal form} of Theorem~\ref{thm:first-routed-return-normal-form}, then \(w\) remains screening-only and cannot be carrier-upgraded. In particular, no carrier-upgraded routed-balance witness exists on the carrier-stable routed-return branch.
\end{proposition}

\begin{proof}
By Theorem~\ref{thm:first-routed-return-normal-form}, the carrier-stable routed return normal form means precisely that there exists one common closure-bearing return-side carrier through which the routed realization is read and that every admitted visible routed mismatch collapses beneath that carrier to the carrier-stable visible normal form. Apply this to the routed-balance witness \(w\). Theorem~\ref{thm:carrier-stable-diagnostic-collapse} then says that, over that same stable carrier package, \(w\) has only one current structural role: it witnesses subordinate screening-loss normal form and does not define an independent carrier-relevant visible species. But Corollary~\ref{cor:carrier-escalation-threshold} identifies carrier upgrade exactly with the first failure of every common closure-bearing return-side carrier-supported realization. Since the carrier-stable routed return normal form explicitly retains such a common carrier package, the threshold for carrier upgrade is not met. Hence \(w\) remains screening-only and cannot be carrier-upgraded. Because \(w\) was arbitrary, no carrier-upgraded routed-balance witness exists on the carrier-stable routed-return branch.
\end{proof}

\begin{corollary}[The live OP~5 Step~(B) burden survives only on the obstructed routed-return branch]
\label{cor:live-op5b-burden-survives-only-on-obstructed-routed-return-branch}
Under the hypotheses of Proposition~\ref{prop:carrier-stable-routed-returns-prevent-routed-balance-witness-escalation}, a genuinely live OP~5 Step~(B) burden can survive only on the \emph{obstructed routed return normal form} of Theorem~\ref{thm:first-routed-return-normal-form}. Equivalently, after the gates \(\mathrm{TR}_4\Rightarrow\mathrm{BS}_1\Rightarrow\mathrm{BS}_3\Rightarrow\mathrm{BL}_2\), any routed-balance witness that still honestly threatens carrier/package loss must already belong to the negative routed-return branch; and when that negative branch localizes at the readout/kernel stratum, \(\mathrm{BL}_5\) compresses the remaining burden to absence of one common shell-stable carrier package.
\end{corollary}

\begin{proof}
By Proposition~\ref{prop:live-op5b-kernel-carrier-upgraded-witness-problem}, the live OP~5 Step~(B) kernel is exactly the production or exclusion of a carrier-upgraded routed-balance witness. Proposition~\ref{prop:carrier-stable-routed-returns-prevent-routed-balance-witness-escalation} excludes such a witness on the carrier-stable routed-return branch. Therefore any genuinely live routed-balance burden must sit on the only remaining routed-return branch, namely the obstructed routed return normal form of Theorem~\ref{thm:first-routed-return-normal-form}. Corollary~\ref{cor:first-normalized-reopening-route-live-op5b-kernel} already records that, after the masked-single-carrier gate \(\mathrm{BL}_2\), the next honest localization step is \(\mathrm{BL}_4\), and that a readout/kernel localization is then compressed by \(\mathrm{BL}_5\) to absence of one common shell-stable carrier package. This is exactly the stated reformulation.
\end{proof}

\begin{remark}[What this second OP~5 Step~(B) sharpening buys]
\label{rem:what-second-op5b-sharpening-buys}
The live OP~5 Step~(B) burden is now smaller than in r05. It is no longer ``any routed-balance witness that might still upgrade.'' It is ``a routed-balance witness on the obstructed routed-return branch that survives the masked-single-carrier test and, if needed, localizes further.'' This is also the point at which the Monte-Carlo/fusion probes materially accelerate the search: once the probes classify a run as carrier-stable on the routed-return side, one no longer needs to pursue carrier-grade OP~5 escalation there. The probes still do not prove the theorematic branch separation by themselves, but they do cut away large amounts of blind BS/BL reopening by steering attention directly to the small class of routed returns that still look obstruction-like after the carrier-stable branch has been excluded.
\end{remark}



\noindent\textbf{Reader cue.} Read the next block as the first localization split of the remaining OP~5 Step~(B) burden on the negative routed-return side. Its purpose is not yet to declare the obstructed routed return harmless, but to separate the part that is already inherited from the upstream OP3/QAM obstruction package from the part that would remain as a genuinely downstream carrier-grade OP~5 obstruction.

\subsection*{19AZ$^{(6c)}$. Obstructed routed returns split into upstream-screened versus downstream carrier-grade OP~5 Step~(B) witnesses}
\label{sec:obstructed-routed-returns-split-upstream-screened-vs-downstream-carrier-grade-op5b}
After 19AZ$^{(6b)}$, the live OP~5 Step~(B) burden survives only on the obstructed routed-return branch. But even this negative branch is still too coarse to be the final witness kernel. The already extracted global obstruction package says that, once the admissible run is inside the bundled OP1$\to$branch$\to$OP3/QAM$\to$OP5 channel, the first licensed failure is already localized either upstream in the OP3/QAM obstruction package or downstream at the carrier-upgraded OP5 obstruction layer. Therefore the obstructed routed-return branch must itself be split: either the routed obstruction is only the return-side shadow of an upstream OP3/QAM obstruction already fixed before carrier-grade OP5 escalation, or it is a genuinely downstream carrier-grade OP5 obstruction that survives that upstream screening.

\begin{definition}[Upstream-screened versus downstream carrier-grade obstructed routed-return witness]
\label{def:upstream-screened-vs-downstream-carrier-grade-obstructed-routed-return-witness}
Assume the hypotheses of Corollary~\ref{cor:live-op5b-burden-survives-only-on-obstructed-routed-return-branch}. Let $(b,r,s;w)$ be a routed branch-shell realization on the \emph{obstructed routed return normal form} together with a routed-balance witness \(w\).
\begin{enumerate}[leftmargin=2em,label=(\roman*)]
\item We say that $(b,r,s;w)$ is an \emph{upstream-screened obstructed routed-return witness} if the currently licensed negative routed-return reading of $(b,r,s;w)$ is already forced by the upstream OP3/QAM obstruction package of Proposition~\ref{prop:global-obstruction-tool}, so that the witness does not yet certify an independent downstream carrier-grade OP~5 obstruction.
\item We say that $(b,r,s;w)$ is a \emph{downstream carrier-grade obstructed routed-return witness} if the witness remains honestly live after the upstream OP3/QAM obstruction screening and therefore certifies a genuinely downstream carrier-upgraded OP~5 obstruction layer in the sense of Proposition~\ref{prop:global-obstruction-tool}.
\end{enumerate}
\end{definition}

\begin{proposition}[The remaining obstructed routed-return burden is an upstream/downstream split]
\label{prop:remaining-obstructed-routed-return-burden-is-upstream-downstream-split}
Assume Proposition~\ref{prop:live-op5b-kernel-carrier-upgraded-witness-problem}, Corollary~\ref{cor:live-op5b-burden-survives-only-on-obstructed-routed-return-branch}, Proposition~\ref{prop:global-obstruction-tool}, and Proposition~\ref{prop:global-transfer-tool}. Let $(b,r,s;w)$ be a genuinely live routed-balance witness on the obstructed routed return normal form. Then exactly one of the following two cases is currently licensed.
\begin{enumerate}[leftmargin=2em,label=(\roman*)]
\item \textbf{Upstream-screened case.} The negative routed-return reading of $(b,r,s;w)$ is already accounted for by the upstream OP3/QAM obstruction package, and the witness is therefore not yet an independent downstream carrier-grade OP~5 obstruction.
\item \textbf{Downstream carrier-grade case.} The witness survives the upstream OP3/QAM obstruction screening and therefore remains as a genuinely downstream carrier-grade OP~5 obstruction witness.
\end{enumerate}
No third autonomous obstructed routed-return witness class is currently licensed.
\end{proposition}

\begin{proof}
By Corollary~\ref{cor:live-op5b-burden-survives-only-on-obstructed-routed-return-branch}, every genuinely live OP~5 Step~(B) witness must already lie on the obstructed routed return normal form. Once the run is inside the bundled OP1$\to$branch$\to$OP3/QAM$\to$OP5 channel, Proposition~\ref{prop:global-obstruction-tool} says that the first licensed failure is already localized either upstream in the OP3/QAM obstruction package or downstream at the carrier-upgraded OP5 obstruction layer. Proposition~\ref{prop:global-transfer-tool} adds that no third autonomous terminal visible class is licensed outside that bundled obstruction reading. Applying this dichotomy to the obstructed routed-return witness $(b,r,s;w)$ yields exactly the two stated cases. The first is the upstream-screened case, where the negative routed-return picture is merely the return-side shadow of an obstruction already fixed upstream. The second is the downstream carrier-grade case, where the witness remains live after that upstream screening and therefore truly belongs to the OP~5 obstruction layer. Exclusivity and exhaustiveness are precisely the content of the bundled global obstruction dichotomy.
\end{proof}

\begin{corollary}[The genuinely new OP~5 Step~(B) burden survives only on the downstream branch]
\label{cor:genuinely-new-op5b-burden-survives-only-on-downstream-branch}
Under the hypotheses of Proposition~\ref{prop:remaining-obstructed-routed-return-burden-is-upstream-downstream-split}, any genuinely \emph{new} OP~5 Step~(B) burden survives only in the downstream carrier-grade case of Definition~\ref{def:upstream-screened-vs-downstream-carrier-grade-obstructed-routed-return-witness}. Equivalently, once the obstructed routed return has been shown not to be merely the return-side shadow of the upstream OP3/QAM obstruction package, the remaining burden is already localized to the downstream carrier-upgraded OP~5 obstruction layer.
\end{corollary}

\begin{proof}
In the upstream-screened case, the negative routed-return witness is already accounted for by the upstream OP3/QAM obstruction package and therefore does not constitute a genuinely new downstream OP~5 burden. Hence only the downstream carrier-grade case can carry a genuinely new OP~5 Step~(B) burden after the upstream screening has been discharged. This is exactly the stated reformulation.
\end{proof}

\begin{remark}[What this third OP~5 Step~(B) sharpening buys]
\label{rem:what-third-op5b-sharpening-buys}
The live OP~5 Step~(B) burden is now smaller again. It is no longer merely ``a witness on the obstructed routed-return branch.'' It is ``a witness on that negative branch that is not already explained by the upstream OP3/QAM obstruction package.'' This is also where the Monte-Carlo/fusion probes become strategically useful on the negative side: they do not prove the upstream/downstream split, but they help tell much earlier whether a run looks like an inherited OP3/QAM obstruction shadow or like a genuinely downstream carrier-grade obstruction. In practical search terms, this means that many apparently live negative routed returns can now be discarded upstream before one spends time reopening the deepest OP~5 carrier-grade layer.
\end{remark}


\noindent\textbf{Reader cue.} Read the next block as the discharge of the upstream-screened side of the remaining OP~5 Step~(B) burden. Its purpose is not yet to close the downstream carrier-grade obstruction kernel, but to remove every negative routed-return witness whose content is already absorbed by the upstream OP3/QAM obstruction package.

\subsection*{19AZ$^{(6d)}$. Upstream-screened obstructed routed returns are discharged before the downstream OP~5 Step~(B) carrier-grade kernel begins}
\label{sec:upstream-screened-obstructed-routed-returns-discharged-before-downstream-op5b-kernel}
After 19AZ$^{(6c)}$, the live negative routed-return burden has only two currently licensed readings: upstream-screened or downstream carrier-grade. But the upstream-screened case is not merely weaker; it is not an independent OP~5 Step~(B) burden at all. By definition it is already the routed return shadow of an obstruction whose first licensed failure sits upstream in the OP3/QAM package. Therefore the honest live OP~5 Step~(B) kernel starts only after that upstream explanation has been discharged.

\begin{proposition}[Upstream-screened obstructed routed-return witnesses do not define an independent OP~5 Step~(B) burden]
\label{prop:upstream-screened-obstructed-routed-return-witnesses-do-not-define-independent-op5b-burden}
Assume Proposition~\ref{prop:remaining-obstructed-routed-return-burden-is-upstream-downstream-split}, Corollary~\ref{cor:genuinely-new-op5b-burden-survives-only-on-downstream-branch}, Proposition~\ref{prop:global-obstruction-tool}, and Proposition~\ref{prop:global-transfer-tool}. Let $(b,r,s;w)$ be an upstream-screened obstructed routed-return witness in the sense of Definition~\ref{def:upstream-screened-vs-downstream-carrier-grade-obstructed-routed-return-witness}. Then $(b,r,s;w)$ does not define an independent live OP~5 Step~(B) burden; its entire negative routed-return content is already discharged by the upstream OP3/QAM obstruction package.
\end{proposition}

\begin{proof}
By Definition~\ref{def:upstream-screened-vs-downstream-carrier-grade-obstructed-routed-return-witness}, an upstream-screened obstructed routed-return witness is precisely one whose currently licensed negative routed-return reading is already forced by the upstream OP3/QAM obstruction package. Proposition~\ref{prop:global-obstruction-tool} says that once the admissible run is inside the bundled OP1$\to$branch$\to$OP3/QAM$\to$OP5 channel, the first licensed failure is already localized either upstream in that OP3/QAM obstruction package or downstream at the carrier-upgraded OP5 obstruction layer. Since the witness is, by hypothesis, upstream-screened, the first licensed failure is already on the upstream side of that dichotomy. Proposition~\ref{prop:global-transfer-tool} excludes any third autonomous terminal visible class outside that bundled obstruction reading. Therefore no additional independent OP~5 Step~(B) burden remains attached to $(b,r,s;w)$ once the upstream package has been acknowledged. In other words, the routed-return shadow carried by $(b,r,s;w)$ is fully discharged upstream and does not itself define a genuinely new OP~5 Step~(B) burden.
\end{proof}

\begin{corollary}[The live OP~5 Step~(B) kernel reduces to the downstream carrier-grade obstructed routed-return branch]
\label{cor:live-op5b-kernel-reduces-to-downstream-carrier-grade-obstructed-routed-return-branch}
Under the hypotheses of Proposition~\ref{prop:upstream-screened-obstructed-routed-return-witnesses-do-not-define-independent-op5b-burden}, the honest live OP~5 Step~(B) kernel is exactly the production or exclusion of a downstream carrier-grade obstructed routed-return witness. Equivalently, once the upstream-screened negative routed-return cases have been discharged, the only remaining live OP~5 Step~(B) burden sits on the downstream carrier-grade branch isolated in Definition~\ref{def:upstream-screened-vs-downstream-carrier-grade-obstructed-routed-return-witness}.
\end{corollary}

\begin{proof}
By Proposition~\ref{prop:remaining-obstructed-routed-return-burden-is-upstream-downstream-split}, every genuinely live obstructed routed-return witness lies in exactly one of two currently licensed classes: upstream-screened or downstream carrier-grade. Proposition~\ref{prop:upstream-screened-obstructed-routed-return-witnesses-do-not-define-independent-op5b-burden} removes the first class as an independent OP~5 Step~(B) burden. Therefore only the downstream carrier-grade class remains available as the honest live OP~5 Step~(B) kernel. This is exactly the stated reformulation.
\end{proof}

\begin{remark}[What this fourth OP~5 Step~(B) sharpening buys]
\label{rem:what-fourth-op5b-sharpening-buys}
The live OP~5 Step~(B) burden is now concentrated on one branch only. It is no longer ``an obstructed routed-return witness that is not yet known to be harmless.'' It is ``a downstream carrier-grade obstructed routed-return witness after the upstream OP3/QAM shadow cases have been discharged.'' This is the first point on the OP~5 side where the negative branch has been reduced to a single honest carrier-grade kernel. Strategically, this is also where the Monte-Carlo/fusion probes become most valuable as a search accelerator: once a negative routed return looks upstream-screened, it can be discarded without reopening the downstream carrier-grade layer, so the expensive search is reserved only for the small class of runs that still look honestly downstream after the upstream package has been factored out.
\end{remark}



\noindent\textbf{Reader cue.} Read the next block as the compression of the remaining downstream OP~5 Step~(B) kernel to one shell-package deficit question on the retained branch-shell corridor. Its purpose is not yet to declare OP~5 Step~(B) globally closed, but to show that once the upstream shadow cases are gone, the honest remaining burden is exactly the absence of one common shell-stable carrier package.

\subsection*{19AZ$^{(6e)}$. Downstream carrier-grade obstructed routed returns compress to one common shell-stable carrier-package deficit on the retained branch-shell corridor}
\label{sec:downstream-carrier-grade-obstructed-routed-returns-compress-shell-stable-carrier-package-deficit}
After 19AZ$^{(6d)}$, the live OP~5 Step~(B) burden is concentrated on the downstream carrier-grade obstructed routed-return branch. The smallest honest next question is therefore no longer whether some vague routed mismatch survives, but whether the retained branch-shell corridor still admits one common shell-stable carrier package once that downstream witness has survived every earlier screening. On the current line this is exactly the readout/kernel compression question behind \(\mathrm{BL}_5\): if the branch and shell strata themselves remain intact on the retained routed corridor, then every remaining downstream carrier-grade witness is nothing but failure of one common shell-stable carrier package.

\begin{definition}[Downstream shell-package deficit witness for the live OP~5 Step~(B) kernel]
\label{def:downstream-shell-package-deficit-witness-live-op5b-kernel}
Fix the live OP~5 Step~(B) setup of Definition~\ref{def:routed-balance-witness-live-op5b-kernel}. A quadruple $(b,r,s;w)$ is called a \emph{downstream shell-package deficit witness for the live OP~5 Step~(B) kernel} if
\begin{enumerate}[leftmargin=2em,label=(\roman*)]
  \item $(b,r,s;w)$ is a downstream carrier-grade obstructed routed-return witness in the sense of Definition~\ref{def:upstream-screened-vs-downstream-carrier-grade-obstructed-routed-return-witness}, and
  \item on the retained routed branch-shell corridor determined after the gates \(\mathrm{TR}_4\Rightarrow\mathrm{BS}_1\Rightarrow\mathrm{BS}_3\Rightarrow\mathrm{BL}_2\), the branch and shell strata remain intact, so the remaining localized burden is entirely in the readout/kernel stratum and therefore equals failure of one common shell-stable carrier package in the sense of Lemma~\ref{lem:qam-readout-kernel-obstruction-compression}.
\end{enumerate}
\end{definition}

\begin{proposition}[Downstream carrier-grade OP~5 witnesses compress to one common shell-stable carrier-package deficit on the retained branch-shell corridor]
\label{prop:downstream-carrier-grade-op5-witnesses-compress-to-one-common-shell-stable-carrier-package-deficit}
Assume Corollary~\ref{cor:live-op5b-kernel-reduces-to-downstream-carrier-grade-obstructed-routed-return-branch}, Corollary~\ref{cor:first-normalized-reopening-route-live-op5b-kernel}, Theorem~\ref{thm:first-routed-return-normal-form}, Lemma~\ref{lem:qam-shell-obstruction-stratification}, and Lemma~\ref{lem:qam-readout-kernel-obstruction-compression}. Let $(b,r,s;w)$ be a genuinely live downstream carrier-grade obstructed routed-return witness for the OP~5 Step~(B) kernel. If, on the retained routed branch-shell corridor carried through \(\mathrm{TR}_4\Rightarrow\mathrm{BS}_1\Rightarrow\mathrm{BS}_3\Rightarrow\mathrm{BL}_2\), the branch and shell strata remain intact, then $(b,r,s;w)$ is exactly a downstream shell-package deficit witness in the sense of Definition~\ref{def:downstream-shell-package-deficit-witness-live-op5b-kernel}. Equivalently, under the retained branch-shell intactness assumptions, the whole remaining downstream OP~5 Step~(B) burden compresses to absence of one common shell-stable carrier package.
\end{proposition}

\begin{proof}
By Corollary~\ref{cor:live-op5b-kernel-reduces-to-downstream-carrier-grade-obstructed-routed-return-branch}, every genuinely live OP~5 Step~(B) burden is already a downstream carrier-grade obstructed routed-return witness. Corollary~\ref{cor:first-normalized-reopening-route-live-op5b-kernel} records that after the retained gates \(\mathrm{TR}_4\Rightarrow\mathrm{BS}_1\Rightarrow\mathrm{BS}_3\Rightarrow\mathrm{BL}_2\), the first further honest localization step is \(\mathrm{BL}_4\), and that whenever the first remaining localized obstruction lies in the readout/kernel stratum, the compression gate \(\mathrm{BL}_5\) rewrites the burden as absence of one common shell-stable carrier package. Now assume that on the retained routed branch-shell corridor the branch and shell strata themselves remain intact. Then Lemma~\ref{lem:qam-shell-obstruction-stratification} forces every remaining first local obstruction into the readout/kernel stratum, and Lemma~\ref{lem:qam-readout-kernel-obstruction-compression} identifies that stratum exactly with failure of one common shell-stable carrier package. Therefore the surviving downstream carrier-grade witness $(b,r,s;w)$ is not a broader residual species: on the retained branch-shell corridor it is precisely a downstream shell-package deficit witness. This is the stated compression.
\end{proof}

\begin{corollary}[The live OP~5 Step~(B) kernel reduces to one common shell-stable carrier-package deficit problem on the retained branch-shell corridor]
\label{cor:live-op5b-kernel-reduces-to-one-common-shell-stable-carrier-package-deficit-problem}
Under the hypotheses of Proposition~\ref{prop:downstream-carrier-grade-op5-witnesses-compress-to-one-common-shell-stable-carrier-package-deficit}, the honest live OP~5 Step~(B) kernel on the retained branch-shell corridor is exactly the production or exclusion of a downstream shell-package deficit witness. Equivalently, once the upstream-screened cases have been discharged and the branch/shell strata remain intact, OP~5 Step~(B) is reduced to the single question whether one common shell-stable carrier package is absent.
\end{corollary}

\begin{proof}
Corollary~\ref{cor:live-op5b-kernel-reduces-to-downstream-carrier-grade-obstructed-routed-return-branch} says that the live OP~5 Step~(B) burden is already concentrated on downstream carrier-grade obstructed routed-return witnesses. Proposition~\ref{prop:downstream-carrier-grade-op5-witnesses-compress-to-one-common-shell-stable-carrier-package-deficit} then shows that, under retained branch-shell intactness, every such witness is exactly a downstream shell-package deficit witness. Hence the remaining kernel is neither a diffuse routed-return problem nor a broader carrier-grade witness search; it is exactly the stated one-package deficit problem.
\end{proof}

\begin{remark}[What this fifth OP~5 Step~(B) sharpening buys]
\label{rem:what-fifth-op5b-sharpening-buys}
This is the first point on the OP~5 side where the surviving kernel is no longer merely a branch classification problem. It has become a concrete deficit question. The live burden is now: not ``some downstream obstruction somewhere on the negative routed-return branch,'' but ``the absence of one common shell-stable carrier package on the retained branch-shell corridor.'' This is also why the Monte-Carlo/fusion probes feel so effective here. Once the probes indicate that the run is both downstream and branch/shell-intact, the search no longer needs to explore broad alternative witness pictures; it can go straight to the small shell-package deficit kernel that \(\mathrm{BL}_5\) was designed to expose.
\end{remark}




oindent\textbf{Reader cue.} Read the next block as the final conservative test of the live OP~5 Step~(B) kernel on the retained branch-shell corridor. Its purpose is not to claim a new global closure principle from scratch, but to show that once the surviving downstream kernel is already compressed to one shell-package deficit question, the inherited OP~3/QAM shell-closure package removes exactly that final deficit under the retained shell-transfer discipline.

\subsection*{19AZ$^{(6f)}$. The inherited shell-closure package discharges the downstream shell-package deficit on the retained branch-shell corridor}
\label{sec:inherited-shell-closure-package-discharges-downstream-shell-package-deficit}
After 19AZ$^{(6e)}$, the live OP~5 Step~(B) kernel has been compressed as far as possible without reusing the earlier OP~3/QAM shell package. The honest remaining question is therefore very narrow: can the retained branch-shell corridor both inherit the shell-closure package and yet still fail to admit one common shell-stable carrier package? On the present line the answer is no. Once the retained branch-shell corridor is genuinely read through the already established shell certificate / shell-stable kernel-readout package, the missing common shell-stable carrier of 19AZ$^{(6e)}$ is exactly the kind of deficit that the inherited package forbids.

\begin{definition}[Retained shell-transfer discipline for the live OP~5 Step~(B) kernel]
\label{def:retained-shell-transfer-discipline-live-op5b-kernel}
Fix the live OP~5 Step~(B) setup of Definition~\ref{def:routed-balance-witness-live-op5b-kernel}. We say that the retained routed branch-shell corridor satisfies the \emph{retained shell-transfer discipline} if, after the retained gates \(\mathrm{TR}_4\Rightarrow\mathrm{BS}_1\Rightarrow\mathrm{BS}_3\Rightarrow\mathrm{BL}_2\), the routed branch-shell realization inherits the shell admissibility certificate and the equivalent shell-stable kernel/readout package from Theorem~\ref{thm:qam-op3-shell-closure}. Equivalently, on that retained corridor there exists one common shell-stable kernel/readout carrier through which all closure-relevant routed readout factors.
\end{definition}

\begin{proposition}[The inherited shell-closure package excludes downstream shell-package deficit witnesses on the retained branch-shell corridor]
\label{prop:inherited-shell-closure-package-excludes-downstream-shell-package-deficit-witnesses}
Assume Corollary~\ref{cor:live-op5b-kernel-reduces-to-one-common-shell-stable-carrier-package-deficit-problem}, Definition~\ref{def:retained-shell-transfer-discipline-live-op5b-kernel}, Theorem~\ref{thm:qam-op3-shell-closure}, and Theorem~\ref{thm:qam-shell-iff-package}. Then no downstream shell-package deficit witness in the sense of Definition~\ref{def:downstream-shell-package-deficit-witness-live-op5b-kernel} can exist on the retained branch-shell corridor. Equivalently, once the retained branch-shell corridor inherits the shell certificate / shell-stable kernel-readout package, the remaining OP~5 Step~(B) kernel from 19AZ$^{(6e)}$ is discharged.
\end{proposition}

\begin{proof}
Assume toward contradiction that a downstream shell-package deficit witness exists on the retained branch-shell corridor. By Definition~\ref{def:downstream-shell-package-deficit-witness-live-op5b-kernel}, this means that the surviving downstream carrier-grade routed-return witness is already compressed to absence of one common shell-stable carrier package on that corridor. But by Definition~\ref{def:retained-shell-transfer-discipline-live-op5b-kernel}, the same retained corridor inherits the shell admissibility certificate and hence, by Theorem~\ref{thm:qam-op3-shell-closure} together with Theorem~\ref{thm:qam-shell-iff-package}, admits one common shell-stable kernel/readout carrier through which the closure-relevant routed readout factors. This is exactly the negation of the shell-package deficit asserted by the witness. The contradiction shows that no downstream shell-package deficit witness can exist once the retained shell-transfer discipline holds. Therefore the remaining OP~5 Step~(B) kernel from Corollary~\ref{cor:live-op5b-kernel-reduces-to-one-common-shell-stable-carrier-package-deficit-problem} is discharged on that corridor.
\end{proof}

\begin{corollary}[Closure of the live OP~5 Step~(B) kernel under the standing retained branch-shell and shell-transfer disciplines]
\label{cor:closure-live-op5b-kernel-under-standing-retained-branch-shell-and-shell-transfer-disciplines}
Assume the standing retained branch-shell intactness hypotheses from Proposition~\ref{prop:downstream-carrier-grade-op5-witnesses-compress-to-one-common-shell-stable-carrier-package-deficit} together with the retained shell-transfer discipline of Definition~\ref{def:retained-shell-transfer-discipline-live-op5b-kernel}. Then the live OP~5 Step~(B) kernel closes:
\[
\Delta_{\Theta R}=0.
\]
Equivalently, under the current retained no-upstream-shadow, branch/shell-intact, and inherited shell-package disciplines, no honest live OP~5 Step~(B) witness survives.
\end{corollary}

\begin{proof}
By Corollary~\ref{cor:live-op5b-kernel-reduces-to-one-common-shell-stable-carrier-package-deficit-problem}, the live OP~5 Step~(B) kernel under the retained branch-shell hypotheses is exactly the production or exclusion of a downstream shell-package deficit witness. Proposition~\ref{prop:inherited-shell-closure-package-excludes-downstream-shell-package-deficit-witnesses} excludes precisely that witness class once the retained shell-transfer discipline holds. Hence no live OP~5 Step~(B) witness remains, so the kernel closes and \(\Delta_{\Theta R}=0\).
\end{proof}

\begin{remark}[What this sixth OP~5 Step~(B) sharpening buys]
\label{rem:what-sixth-op5b-sharpening-buys}
This is the OP~5 analogue of the OP~1 local closure step. The result is intentionally conservative. It does not say that every historical reading of OP~5 is now globally absorbed into one theorematic sentence. It says something narrower and structurally cleaner: once the live OP~5 Step~(B) burden has been reduced to one downstream shell-package deficit on the retained branch-shell corridor, the inherited OP~3/QAM shell-closure package forbids exactly that final deficit. Strategically, this is also the point where the Monte-Carlo/fusion probes show their full value on the negative routed-return side. They do not prove the inherited shell package, but they steer the search quickly onto the tiny downstream/branch-shell-intact corridor where the shell-transfer discipline can be applied and the final witness class disappears.
\end{remark}
\begin{remark}[Historical transition note for the archived v2.27.x release-control fragment]
Read the next local post-release fragment as archival proof-memory from an earlier public hardening cycle. Its purpose is to preserve release genealogy, checkpoint recall, and local baseline memory, not to function as a current roadmap, freeze checklist, or active control surface for the present v\docversion\ line.
\end{remark}

